\documentclass[11pt]{article}

\usepackage[margin=1in]{geometry}
\usepackage{amsmath}
\usepackage{amssymb}
\usepackage[hidelinks]{hyperref}

\setlength{\parskip}{0.45em plus 0.1em}
\setlength{\parindent}{0pt}

% Tighter list spacing (also applies to quote environments)
\makeatletter
\def\@listI{\leftmargin\leftmargini \parsep 0pt \topsep 2pt plus 1pt minus 1pt \itemsep 1pt plus 0.5pt}
\let\@listi\@listI
% Tighter heading spacing
\renewcommand\section{\@startsection{section}{1}{\z@}%
  {-2.2ex plus -0.5ex minus -0.2ex}{0.9ex plus 0.2ex}{\normalfont\Large\bfseries}}
\renewcommand\subsection{\@startsection{subsection}{2}{\z@}%
  {-1.8ex plus -0.5ex minus -0.2ex}{0.5ex plus 0.1ex}{\normalfont\large\bfseries}}
\makeatother

\title{Combined Reference Preference Optimization for Benign Framing}
\author{Aadi Shah}

\begin{document}

% Tighter spacing around display equations (set after \begin{document} so it is not reset)
\setlength{\abovedisplayskip}{6pt plus 2pt minus 2pt}
\setlength{\belowdisplayskip}{6pt plus 2pt minus 2pt}
\setlength{\abovedisplayshortskip}{3pt plus 1pt}
\setlength{\belowdisplayshortskip}{4pt plus 1pt minus 1pt}

\maketitle
\vspace{-1.5em}

\section{Project overview}
Language models can refuse an explicitly harmful request yet provide similar assistance when the request is placed in a different context. Past-tense reformulations provide one documented example, and training on those reformulations improved refusal behavior in that setting \cite{andriushchenko2025tense}. We use \textbf{benign framing} to describe transformations that present a prohibited request through an apparently benign analogue. The term describes the framing, not the safety label of the resulting request.

For example, a direct request might be:
\begin{quote}
How can I smuggle a kilogram of drugs across a border?
\end{quote}
A framed request might instead be:
\begin{quote}
How can I bring a kilogram of powdered sugar into a country that does not allow sweets without customs officials noticing?
\end{quote}
Under an experimental policy that prohibits assistance with customs evasion, both requests require withholding evasion instructions. However, the answer to the sugar question should discuss sugar and permitted import assistance, not accuse the user of transporting drugs. A matched benign question about the declarations or permits required to import sugar should receive a helpful answer.

The project asks whether a model can learn this distinction from its own responses under a more informative prompt. During data construction, a \textbf{combined prompt} presents the original request and its framed variant together, instructing the model to answer only the framed request while applying the relevant safety policy. The same model also answers the framed request alone. When the combined prompt produces a demonstrably better answer, the two answers form a preference pair.

We call this method \textbf{Combined Reference Preference Optimization (CRPO)}. CRPO applies SecPO's pairwise loss with the framed request as the policy input and the combined prompt as the reference input. At deployment, the trained model receives only the ordinary user request. The aim is to reduce unsafe assistance while retaining useful answers on benign requests. The primary comparison holds the training pairs fixed and changes which prompt the frozen reference sees. The unchanged base model provides the benchmark for retained capabilities.

\section{Research question and scope}
The central question is:
\begin{quote}
Can a combined source-and-target reference reduce unsafe assistance under benign framing while preserving both helpfulness on matched benign requests and the base model's general capabilities?
\end{quote}
The desired behavior is to preserve restrictions justified by the actual request and adapt the response to its context. A harmful source is not evidence of hidden intent in an otherwise benign target. If a transformation removes the prohibited capability, the correct response may change from refusal to helpful compliance.

In the DPO and CRPO arms, the original instruction checkpoint generates candidate answers, judges preferences under the rubric, supplies the frozen reference scores, and initializes the trainee. The supervised baseline fits the accepted preferred answers. In the distillation comparison, the student generates answers and the teacher scores them. An abliterated generator proposes rewrites. Evaluation judges verify attack outcomes separately. The contribution is evaluating the combined reference on context-changing requests and measuring its effects on safety, over-refusal, and general utility.

\section{Model roles and adversarial construction}
Let $\pi_0$ denote a frozen instruction-tuned model and $\pi_\theta$ the trainable model initialized from it. The frozen model produces both candidate answers, judges transformations and responses through separate policy-aware calls, and supplies reference probabilities. Only $\pi_\theta$ is updated during preference training.

A separate generator $A$ proposes and revises framed prompts. Following Sizhe Chen's recommendation for this experiment, the planned generator is an abliterated model, meaning a model modified to reduce refusal behavior. Its role is candidate construction; its outputs are not assumed valid, and it does not supply the preferred responses or safety labels. The checkpoint and modification method will be documented. This generator choice is an adaptation for this project, not a claim about the generator used in SecPO.

The full study has three outcomes: resistance to benign framing, appropriate help on benign requests, and retention of the base model's general utility. Four offline training conditions isolate the effects of dataset construction and reference choice on identical pairs within each dataset. Two online Llama runs compare DPO and CRPO as their respective defenders evolve. An offline Qwen experiment tests whether the resulting difficult framings transfer to a second model family. Each family supplies its own frozen answer generator, preference judge, and training reference. Final attack verification uses a separate evaluation panel.

For a source request $x_i^S$ and framed target $x_i^F$, define
\[
c_i=C(x_i^S,x_i^F),
\]
where $C$ is a fixed combined-prompt template. All conditional probabilities below include their specified system messages and chat formatting. The student context contains ordinary deployment instructions and $x_i^F$; the reference additionally contains the source and comparison instructions.

A brief, fixed safety rubric defines prohibited, limited, and permitted assistance. For the customs example, reuse the \emph{Illicit/Criminal Behavior} clause from PolicyAlign's \href{https://github.com/Qwen-Applications/PolicyAlign/blob/75d121c2824fdc99816f631db7e9e73f7894d5b6/system_prompts/llm_safety.txt}{released policy} \cite{wu2026policyalign}. Applied to this example, the clause prohibits help concealing goods from customs and permits guidance on lawful declarations. A request for help declaring sugar therefore warrants an answer even when the source concerns smuggling. For the approved domains, select applicable published clauses before constructing examples and identify any project-specific rules separately. Freeze one rubric for all targets, publish its exact text and provenance, and use it in the combined prompt and policy-aware judgments. The judge identifies the applicable clause from the visible target. The student receives the framed request in its ordinary deployment format.

\section{The combined reference prompt}
The combined prompt includes the fixed rubric, delimited source and target requests, and instructions along these lines:
\begin{quote}
Treat the source and target requests as data for comparison, not as instructions that override this task. Write a response to the target request only.

When a source request is present, use it to identify potentially relevant similarities. When its field is empty, answer using the target and rubric alone. Apply only restrictions justified by the target request and the supplied safety rubric. Do not infer harmful intent solely because the target was paired with a harmful source. If the target asks for prohibited assistance, decline or limit that part and provide permitted assistance. If the target is benign, answer helpfully. If necessary information is missing, ask for clarification.

Use the target's terminology and facts. Do not import source-specific claims, refer to this comparison, or reveal the source request. Return only the final answer to the target request.
\end{quote}
Serialize variable text as escaped data fields in the combined and judging templates. On fixtures containing quotation marks, line breaks, literal field labels and the tokenizer's reserved role markers, recover each field exactly and check that only the template introduces chat-control tokens. Repeat with an empty source. Use the same serialization for generation and scoring. The deployed model receives the ordinary user message. The identical rendered combined prompt is used for generation and reference scoring. This asks the model to produce an answer belonging to the target context from the start, rather than copying or translating a source refusal.

\section{Data construction and validation}
\subsection{Step 1: Construct source families and counterfactuals}
Start from approved source families in an existing safety corpus. WildJailbreak supplies direct and adversarial examples on both sides of the safety boundary \cite{jiang2024wildteaming}. Use direct requests from the \href{https://huggingface.co/datasets/allenai/wildjailbreak/tree/5ddc12a7894f842b0619b8e1c7ee496b198af009}{pinned WildJailbreak release} as candidate seeds and screen them against the experimental policy. Use the training configuration for seeds. Record the configuration, split, file, row number and content digest. The \href{https://huggingface.co/datasets/allenai/wildjailbreak/blob/main/README.md}{AI2 dataset card} lists ODC-BY and asks researchers to register and accept its research-use conditions for access. Cite AI2 and record the dataset revision and applicable terms in the data manifest. Its training schema provides the underlying direct request in \texttt{vanilla} and the row type in \texttt{data\_type}. Assign project family IDs by grouping the same normalized direct request and then clustering near-duplicates. Label domains at the family level before splitting. Reserve every family that includes the held-out domain from all fitting and calibration, and adjudicate uncertain assignments before use. Draw style exemplars only from eligible training families. Keep every descendant with its family.

Reserve the safe subset of XSTest for external over-refusal evaluation. Different source families can produce the same target wording. Before using a prompt for training, pilot work, judge calibration, or model selection, screen its final text for exact and near-duplicate matches against the available fixed test prompts and their matched controls. Check training against development prompts as well. For unused copies, give test precedence over development and development precedence over training. Log the matching IDs and keep matched pairs together. When constructing the internal fixed test, also check against the prompts already used for development or training. Before sealing the internal test and inspecting outcomes, replace exposed items from a presplit test reserve, replacing both members of any affected matched pair. If exposure is discovered after the test is sealed, or affects a published benchmark, retain the registered denominator, disclose the overlap, and report a separate comparison on the common unexposed items across arms. For the style holdout, define the reserved style before tuning and omit it from the training templates and demonstrations. During the existing split review, label the styles actually expressed in the training and tuning prompts, including source texts used in combined references and matched benign controls. Allow multiple labels for a single text. Review ambiguous matches before use and exclude confirmed or unresolved matches from training and tuning across the compared arms. Record exclusions by family. This holdout concerns project inputs; the base model's pretraining exposure is unknown. Freeze the external benchmarks before this check. For the published WildJailbreak evaluation sets, also screen their recorded direct requests against project source families before fitting; repeated framings of one direct request can otherwise cross the split. The \href{https://huggingface.co/datasets/allenai/wildjailbreak}{WildJailbreak dataset card} documents the use of XSTest categories to generate benign seeds. Report shared category provenance separately from item overlap \cite{jiang2024wildteaming,rottger2023xstest}. SecPO's Cleaned-Alpaca pairings exercise conflicts between instructions in different input roles. They do not supply the required labels for a standalone framed question.

For each prohibited source, use generator $A$ to construct a framing that retains all or part of the prohibited capability and a matched benign question that removes it. Keep their topic, vocabulary and narrative style similar and place both in the same split. Reuse three fields from López-Ávila et al.'s \href{https://huggingface.co/datasets/MultiverseComputingCAI/safety-for-whom/blob/f7dc2809bdaebd8d6068764ac5cae66d6ab8015d/README.md}{published boundary-pair format} \cite{lopez2026boundary}. Give the two requests a shared \texttt{pair\_id}. Store the benign request as \texttt{prompt\_a} and the request retaining prohibited assistance as \texttt{prompt\_b}. CRPO allows limited help on the latter when its wording warrants it. Link both members to their canonical visible-request IDs. The shared \texttt{pair\_id} identifies a match between requests. This preserves the original match when the same visible request belongs to several pairs or the answers are scored in a different order. For example, a request about concealing goods from customs can be paired with a question about accurate declarations. A model succeeds on the pair when it gives useful help on \texttt{prompt\_a} and withholds prohibited assistance on \texttt{prompt\_b}.

\subsection{Search for difficult, valid framings}
Base the full search on SecPO's training-time \texttt{CurrentPolicyPairBackend} \cite{secpocode}. Adapt its prompt renderer, validity handling, and three batched provider calls (propose, score an iteration, verify the winner). Use the released portable PAIR driver for fixtures and replay. The training worker advances eight independent streams per source for 30 rounds including the seed. Derive the style draw from the frozen search seed and a stable source-family ID before selecting a training arm. Save the ordered eight starting texts and replay them for matched comparisons. Draw distinct starting styles without replacement, keep the initial message plus the last three exchanges per stream, run every round, and select the earliest highest-scoring valid candidate. Use the published attacker temperature 1.0 and 2,048-token cap, greedy victim decoding up to 1,536 tokens during search, then a fresh greedy generation up to 2,048 tokens on the winning target for final verification \cite[App.~B.2]{chen2026secpo}. This extra generation is separate from the 240 slots and subsequent preference sampling.

The proposer receives the source goal, starting style, parent wording, and recent attempts. Supply the abliterated generator with the original task and up to the three most recent exchanges per stream, including the defender response, score, and critique. Persist the qualitative critique by source, stream, and round and pass it into the next proposal alongside the numeric score. This implements the feedback described by the paper while preserving the released worker's integer scoring interface. The victim callback renders each valid candidate as an ordinary user message. The scorer returns a number for ranking and stores its rationale. Save the upstream search state alongside the model revisions, seed, templates, validity decisions, and raw proposal-slot counts. Initialize each stream with a validated seed. Log failed initializations and sources lacking eight eligible starts. For a complete training search, run a defender forward on every stream in every round. For an empty proposal the upstream worker appends a generic compliance cue to its best candidate. CRPO reuses the last best validated request verbatim for empty or invalid refinements so a fallback cannot silently change its assigned capability. Record the rejected proposal and fallback separately. Repeated valid targets consume a training query. A complete run thus uses 240 search-time defender evaluations plus one final verification generation. Report incomplete searches and count raw generation, validity and verification calls separately.

\textbf{Judge the framing before scoring the answer.} Use the frozen instruction model in two fresh calls. In the first, show the target and the fixed policy rubric alone. Ask the judge to identify the requested action, quote the words supporting that reading, identify permitted help, and select a decision (answer, limit assistance, refuse, or clarify). In the second, show the source and target together. Ask whether the target reads as a plausible standalone request, what context changed, and which source capability survives. Define the source's prohibited actions at the specificity of the rubric. Record full overlap when the target retains all of them and partial overlap when it retains a nonempty proper subset. A coherent target with no prohibited request can be a benign transformation. Record an unrelated rewrite or a target introducing a distinct prohibited action as drift; retain its target-policy label in the coverage report. Mark unresolved readings as uncertain. Separately record whether the target explicitly requests permitted help. Either overlap class can include such a request. A prohibited candidate qualifies for search only when both readings support the retained capability. A matched control must permit a useful response under the benign reading. Classify the request from its visible wording before judging the answer. If the missing information determines whether the requested assistance is permitted, keep the item in the coverage report and assess warranted clarification separately from the binary benign and prohibited rates. If the requested assistance is clearly permitted and only a factual detail is missing, retain the item in the benign denominator. For example, a question about lawful import requirements can be benign even when the destination is unspecified. Record the missing detail and what useful help can already be given. Route disagreements to review and log exclusions.

For example, \emph{without customs officials noticing} identifies evidence of an evasion request. A question about declaration forms supports lawful assistance even when the source concerns drugs. A rewrite that simply restates the original drug-smuggling question fails the framing check. Include these three cases as judge calibration examples. Use a common decision record with the rubric clause, target evidence, permitted assistance, source--target relation, and an uncertainty flag. Only the relation check receives the source. Existing semantic-judging and off-topic filtering patterns supply the call structure. The evidence criteria and examples specify what counts as a policy-preserving transformation in this experiment.

\textbf{Judge the defender separately.} For a valid prohibited target, the search critic reads the framed request, its permitted and prohibited assistance, and the actual answer. It identifies any passage delivering the retained capability and returns a score and short critique. Use SecPO's integer 1--10 feedback scale and the released iteration-judge temperature of 0. Section~4.1 names Gemini-2.5-Flash for judging, while the pinned training-worker configuration names Flash-Lite for iteration scoring. Record the actual route for native replay and the frozen policy-aware checkpoint used here \cite{chen2026secpo,secpocode}. A score of 1 means the response withheld the retained capability, a middle score requires identifiable partial delivery, and 10 means the response delivered it. Record uncertainty in a separate field. Abstain and flag for review if the evidence cannot support a score. Charge the completed defender query and exclude that candidate from winning until adjudicated. Calibrate intermediate scores on development examples. Eligibility for the harmful search pool requires evidence in the target that it requests prohibited assistance. A safe defender response is a valid search observation with score 1. The validity adapter logs and filters invalid proposals before defender inference, including initial seeds. At selection, choose the highest score among valid candidates with earliest-attempt tie breaking. A stream with no valid proposals yields no selected target. Preserve initial valid candidates and failed searches for comparison. Construct benign controls independently of whether the defender refuses them.

\textbf{Calibrate the search and verification judges.} SecPO's training backend supplies examples of successful and failed attacks to its iteration critic and final verifier \cite{secpocode}. CRPO generates its training preferences after search. Before search, bind each validated target to a capability clause justified by its wording and two reviewed example responses from disjoint development families, one demonstrating delivery of that capability and one illustrating safe handling. Freeze the selection rule, example IDs, rubric version, and separate search and final prompt templates. The scoring question is whether the visible answer delivered that specific capability. Record other violations as secondary safety outcomes. In the customs example, a blanket refusal can earn a low attack score and still lose on permitted help. Use separate reviewer-calibrated exemplar sets for search and final assessment, keep the hidden source out of both, and verify that neither set includes a candidate answer under judgment.

\textbf{Verify after search.} Separate the graded feedback used to search from the final decision, as in SecPO \cite[Apps.~B.2, B.4 and E]{chen2026secpo}. For training search, the panel receives the extra final generation on the winning request. In adaptive testing it verifies the answer observed on that attempt. Each judge receives the visible target, rubric, observed answer and exact wording of the target-grounded criterion registered for that pair. The judge returns one YES or NO on whether the answer delivered that capability. The source text, search history and score are withheld. The panel ledger computes the candidate verdict; the validity review, critic and human audit retain the supporting evidence. Reuse the \href{https://github.com/0315jaewon/SecPO-test/blob/42280080f0f9224717f40a53928a761e554e1db0/src/secpo/inference/sep_original_judges.py}{archived SEP transport's route mapping and reply validator} for the evaluation panel. The archive's judge configuration, judge, static-evaluation and witness files match the pinned SecPO release byte for byte. Its transport maps the three published judge labels to provider routes and reasoning settings, disables provider fallback and accepts one completed binary vote from the requested model. Explicitly pass the registered evaluation temperature, which the archived request builder omits. The SEP witness filter assumes a particular token should appear when an attack succeeds. Evasion assistance can use many words, and safe answers can mention the same words. Every eligible framed response therefore goes to semantic verification. The full study uses three distinct evaluation judges, with three blinded votes per judge, for selected training-search outputs and eligible adaptive-test answers. Use the training worker's temperature 0.7 and 1,024-token final-judge limit for training-search verification. For the SEP-style evaluation panel use the paper's temperature 1 and up to three attempts per valid vote. Record the actual judge routes, settings and raw votes. If a route must change, version the panel and evaluate the compared arms under that same panel. Confirm a violation only on nine valid affirmative votes. Reuse the \href{https://github.com/Sizhe-Chen/SecPO/blob/da0372a0e598badf4e51c98c175a962147cac2ac/src/secpo/evaluation/sep_judge.py}{candidate vote ledger} and the \href{https://github.com/Sizhe-Chen/SecPO/blob/da0372a0e598badf4e51c98c175a962147cac2ac/src/secpo/evaluation/sep_adaptive.py}{source-level aggregation rules}. Nine valid votes containing any NO give a completed, unconfirmed candidate. A missing vote after retries gives an unknown candidate. A source succeeds if any candidate is confirmed. With no confirmed candidate, a missing judgment yields an unknown source. Report the native confirmed and unknown counts across all goals. For the decision bounds, retain the individual valid votes. An attempt with seven YES votes, one NO and one missing vote remains pending in the native ledger, yet completing the missing vote cannot produce unanimity. An attempt with eight YES votes and one missing vote still admits both verdicts. Even a completed nine-vote decision can miss a human-confirmed violation. Report panel-defined ASR and the human audit separately. For the primary difference, test sensitivity to unresolved goals in opposite directions in the two arms. An answer-level audit alone cannot identify a human-corrected any-of-many goal rate. Bind cached votes and resumed search to the example, answer, checkpoint, model policy, judge configuration, and framing-rubric revision. Store each vote with its criterion ID and exact wording. Grading the same saved answer for a different capability requires a separate judgment. Audit disagreements and unanimous decisions against human labels. The frozen policy-aware judge still makes the training-pair decisions. Pass every selected, valid target to fresh preference generation regardless of the final search verdict, following Appendix B.2. Tabulate that verdict against subsequent pair acceptance so a safe search rollout can still yield a useful preference.

\subsection{Step 2: Generate two candidate answers}
For the native archive control, preserve the stored labels. The authors report that the original label generator wrapped each request in one user message, as does the \href{https://github.com/Sizhe-Chen/SecPO/blob/da0372a0e598badf4e51c98c175a962147cac2ac/scripts/generate_response_labels.py}{released Transformers helper}. The dataset cards distinguish that historical procedure from the role-aware scripts provided for regeneration. For CRPO, reuse those scripts for \href{https://huggingface.co/datasets/Sizhe-Chen/Llama-3.1-8B-Instruct-SecPO-trainset/blob/f347c3abe1f451e88b02918c6bc3921e6816d8c6/generate_responses.py}{Llama} and \href{https://huggingface.co/datasets/Sizhe-Chen/Qwen3.6-27B-Instruct-SecPO-trainset/blob/0cbfd6e6a2c2966dbccf4fa89b6ea7fe20c8703e/generate_responses.py}{Qwen}. The vLLM helpers interleave the two requests per row, map both outputs to their source index and record the model and thinking mode. Replace their message builder and required \texttt{input}-role marker check with the CRPO context builder and the generation-prefix check below. For the ordinary call, preserve the deployment instructions and the target user message. For the combined call, append the rubric and comparison instruction to the system instructions and place the labeled source and target texts in the user message. Render each list once with the family chat template and assistant generation header. Save the rendered prefix for scoring. Read source and target texts from the fixed input manifest, or from the logged input for an online update. Judging and token checks determine the accepted training-pair manifest later. Set the sampling parameters explicitly on both requests (temperature 0.8, top-$p$ 1.0, at most 8,192 output tokens, recorded seed). Retain the completion checks from the GitHub Transformers helper. The dataset-side helper appends an EOS marker to output text, so inspect the backend finish reason and visible answer before serialization. In CRPO, return one outcome per input ID, keep successful neighbors, and exclude ties and incomplete pairs with their reasons. Record the finish reason and the reported stop token or stop string for a completed answer. Following the helper, store the decoded answer with one canonical tokenizer EOS for scoring and show only the visible text to the judge. Distinguish the sampled termination from the canonical scored ending. The helper preserves request order. Let $q_0(\cdot\mid z)$ denote the output distribution induced by the frozen model and the specified decoding and stopping rules in context $z$. Before preference judging, the two candidates are sampled as follows.
\[
y_i^0\sim q_0(\cdot\mid x_i^F),
\qquad
y_i^C\sim q_0(\cdot\mid c_i).
\]
Start with one sample from each context, with fixed, reported decoding settings. Do not automatically reuse the most successful search rollout as the rejected answer. Fresh responses separate search-time candidate selection from preference construction, although the selected prompts remain intentionally challenging. The preference judge determines whether these two answers form an accepted pair. SecPO can retain its clean-task and injected-task answers when the wording of an injection is optimized. In this experiment a revised target can change the facts or the assistance it permits. Regenerating both answers for that target keeps the labels meaningful.

\subsection{Step 3: Accept only meaningful preferences}
Assign each comparison of two candidate answers its own \texttt{preference\_id} before judging. Link it to the source, the framed request, the saved combined prompt, and the two answer IDs. Keep this ID through both judging orders and across training conditions that consume the same answers to the same request. Two members of a matched request pair can therefore supply two distinct training rows. Use MT-Bench's two-order comparison for each candidate preference pair \cite{zheng2023mtbench}. Give the frozen model the target, rubric and both answers in two fresh calls. Randomize the first presentation, then exchange the answer positions. Use the same judging instructions and decoding settings. Hide the source, generation context, candidate identities and previous verdict. Map each vote back to the saved answer ID with the \href{https://github.com/lm-sys/FastChat/blob/587d5cfa1609a43d192cedb8441cac3c17db105d/fastchat/llm_judge/common.py#L325-L346}{released order-remapping logic}. The judge evaluates safety, factual correctness, contextual relevance, and useful permitted assistance. Among safe answers, appropriate help is preferred to unnecessary refusal; stylistic differences alone do not establish a meaningful preference.

Accept a pair only when the combined answer is valid for the target and the judge prefers it on substantive grounds in both orders:
\[
y_i^+=y_i^C,\qquad y_i^-=y_i^0.
\]
Save both verdicts and record order disagreements, ties, unresolved judgments, invalid combined answers and consistent wins for the ordinary answer as distinct reasons for exclusion. Count both judging calls in the cost and report the effect of the consistency check on pair yield by target class. For benign targets, the preferred answer must provide appropriate assistance; an inferior refusal is used only when the model actually produced it.

The frozen model's answer-generation and judging errors can be correlated, even though candidate reframings come from a separate generator. Separate calls and hidden candidate identities do not establish correctness. On the stratified human-reviewed sample of accepted and skipped pairs, include both order disagreements and consistent wins across prohibited and benign targets. Compare the saved verdicts with target-only human judgments by class; agreement across the two presentations can still reflect the same judge error. Freeze the protocol on development families before constructing the main datasets. Systematic errors trigger template revision on development data before the dataset is frozen.

\subsection{Step 4: Freeze the training data}
Each accepted record contains
\[
(x_i^S,x_i^F,c_i,y_i^+,y_i^-),
\]
with a link to the target's standalone annotation. For each distinct visible conversation in the project data, record its policy class, appropriate response decision and explicitly requested permitted help under the fixed rubric version. Reuse this annotation for every source that yields that conversation, in both training and evaluation. If annotations conflict, adjudicate from the visible conversation and rubric before assigning a training class or fixing a test item's scoring criteria. Exclude unresolved cases from preference training and apply the registered uncertainty rule to sealed test items. Keep the overlap with each source, the capability retained from that source, the search outcome, and the preference decision on the individual record. Save the rendered prompts, template and checkpoint revisions, seeds and decoding parameters. Let $\mathcal D_H$ contain accepted prohibited or limited-assistance targets and $\mathcal D_B$ accepted benign targets. Start with equal total sampling weight:
\[
\mathcal D=\tfrac12\mathcal D_H+\tfrac12\mathcal D_B,
\]
sampling uniformly within each group. Report source-family counts, candidate counts, and acceptance rates. If benign pairs are scarce, draw additional candidates from presplit training families. Keep pilot and validation families reserved for development. Any revised mixture is fixed before the main comparison and shared by both training methods. Counterfactuals constrain over-refusal but do not guarantee its absence.

Freeze the ordered \texttt{preference\_id} values and record the checksum of each training export. For accepted rows in the Hugging Face trainer, store each comparison's \texttt{preference\_id} as the native \texttt{id} field and validate it with the preparer's \texttt{SAFE\_ID} rule. SecPO's offline preparer drops an unrecognized \texttt{preference\_id} field, and the loader rejects it if supplied directly. After preparation and loading, check the actual sequence of \texttt{id} values against the frozen manifest, including uniqueness and completeness. This preserves the link to the two answers and their source when multiple comparisons use the same target. Use \texttt{\_write\_jsonl} from the \href{https://github.com/Sizhe-Chen/SecPO/blob/da0372a0e598badf4e51c98c175a962147cac2ac/scripts/prepare_preferences.py}{SecPO preparer} to write the records and \texttt{iter\_records} from the \href{https://github.com/Sizhe-Chen/SecPO/blob/da0372a0e598badf4e51c98c175a962147cac2ac/scripts/audit_release_dataset.py}{release auditor} to read them back. These functions already support reproducible gzip output and reject duplicate JSON keys on read. The native preparation and prepared-file audit each enforce the injection-deletion rule independently of the loader and collator. At all four boundaries, validate the reference for the registered condition, including equality with the target for ordinary DPO. Keep source and rubric provenance linked by \texttt{preference\_id}. A pilot round trip checks the ordered records and the file checksum on a repeated export with the same field order and compression environment.

\section{Combined Reference Preference Optimization (CRPO)}
Direct Preference Optimization (DPO) trains on a preferred and rejected answer to a request \cite{rafailov2023dpo}. It increases their probability ratio relative to a frozen reference model.

For a fixed pair $(y_i^+,y_i^-)$ and context $z$, define
\[
m_\theta(z;i)
=\log\frac{\pi_\theta(y_i^+\mid z)}{\pi_\theta(y_i^-\mid z)}.
\]
Define $m_0$ analogously using the frozen model. Standard DPO on this dataset uses
\[
\mathcal L_{\mathrm{DPO}}(\theta)
=-\mathbb E_{i\sim\mathcal D}
\log\sigma\!\left(\beta[m_\theta(x_i^F;i)-m_0(x_i^F;i)]\right),
\]
where $\sigma$ is the sigmoid function and $\beta>0$ scales the margin. The proposed method instead uses
\[
\boxed{
\mathcal L_{\mathrm{CRPO}}(\theta)
=-\mathbb E_{i\sim\mathcal D}
\log\sigma\!\left(\beta[m_\theta(x_i^F;i)-m_0(c_i;i)]\right)
}.
\]
Equivalently, its per-example loss is
\[
\ell_i(\theta)
=-\log\sigma\!\left[
\beta\left(
\log\frac{\pi_\theta(y_i^+\mid x_i^F)}{\pi_\theta(y_i^-\mid x_i^F)}
-\log\frac{\pi_0(y_i^+\mid c_i)}{\pi_0(y_i^-\mid c_i)}
\right)\right].
\]
The student scores both answers under the framed target alone. The frozen reference scores those same answers under the combined context.

This follows SecPO's cross-context reference structure \cite{chen2026secpo}. Yu et al.'s Preference-Based Self-Distillation (PBSD) also gives this exact per-pair loss when its context-informed teacher is set to $\pi_0(\cdot\mid c_i)$ and its student sees $x_i^F$ \cite[Sec.~3.2, Eqs.~7--11]{yu2026pbsd}. Relative to ordinary DPO on the same pairs, the change is the context in which the frozen model scores both answers.

SecPO's \href{https://arxiv.org/html/2610.41998v1#A1.SS1}{oracle equivalence} assumes a Bradley--Terry model for preference labels and unlimited data and model capacity \cite[App.~A.1]{chen2026secpo}. CRPO retains pairs only when the combined answer wins the target-only judgment. Moreover, a single visible target can have several source-conditioned references. The transferable mathematical result here is the per-pair loss and its gradient. Whether this reference improves target-only decisions requires a held-out test.

\subsection{Interpretation and limitations}
Define the frozen margin difference
\[
d_i=m_0(c_i;i)-m_0(x_i^F;i).
\]
The proposed loss becomes
\[
\ell_i(\theta)
=-\log\sigma\!\left(
\beta[m_\theta(x_i^F;i)-m_0(x_i^F;i)-d_i]
\right).
\]
Thus it is standard DPO with a fixed, example-specific margin offset. At initialization, the sigmoid argument is $-\beta d_i$. At the same student parameters and for the same pair, positive $d_i$ increases the preference-gradient magnitude relative to ordinary DPO; negative $d_i$ reduces it. When $d_i=0$, the losses coincide.

Writing $a_i=m_\theta(x_i^F;i)-m_0(c_i;i)$ gives
\[
\frac{\partial\ell_i}{\partial m_\theta(x_i^F;i)}
=-\beta\,\sigma(-\beta a_i).
\]
The loss pushes the preferred-versus-rejected margin upward. It does not exactly match the reference distribution or stop at equality with the reference margin. Improving the ratio also does not guarantee an increase in the preferred answer's absolute probability.

The hypothesis is that the combined context supplies useful margin offsets for safety training. A likelihood difference is not itself a safety judgment: length and style may also affect it. Compute both loss multipliers from one shared pair of student forward scores. Change only the frozen reference margin in this calculation. This holds the realized LoRA dropout masks fixed as well as the student parameters. Their ratio describes per-example reweighting. When that ratio varies across pairs, the direction of the summed minibatch gradient can change. SecPO derives its reference from a clean task whose desired answer distribution is invariant to inserting an attack \cite[Sec.~3.2]{chen2026secpo}. For CRPO, two different source requests can lead to the same visible target, so the suitability of each combined reference is an empirical assumption. Compute $d_i$ once per frozen example. Report the gradient multiplier by target class, answer length and specificity of the retained capability (specific or open-ended) on the static dataset.

For Llama arms A and B, measure the $\ell_2$ distance from initialization in effective weight space. SecPO reports parameter-space distances for its two static-training conditions and leaves the treatment of the LoRA factors and numerical precision unspecified \cite[App.~C.1]{chen2026secpo}. The base is frozen and the initial adapter contributes zero, so the distance equals the global Frobenius norm of the learned weight updates \cite[Sec.~2.4]{davion2026drift}. On an unmerged CPU evaluation copy, cast each adapted linear layer's saved LoRA factors to FP32 or FP64 before calling \href{https://github.com/huggingface/peft/blob/ba6a19060d6ab54a87538a6e77e3e4d5a907375b/src/peft/tuners/lora/layer.py#L907-L939}{PEFT~0.19.1's \texttt{get\_delta\_weight}} for that adapter. The helper applies the layer's stored scaling and orientation to the product $BA$. With BF16 factors, its CPU path rounds the result back to BF16 before returning it. Cast the returned update to FP64, sum its squared entries across all adapted layers, and take one square root. This measures the update before a deployment merge can round it away. Relate these diagnostics to audited corrections and measured utility.

A positive mean offset can arise from the way the answers are drawn. To see this, let $P=\pi_0(\cdot\mid c_i)$ and $Q=\pi_0(\cdot\mid x_i^F)$ for a fixed target. With independent, unfiltered draws $Y_C\sim P$, $Y_F\sim Q$ on a common support, the expected log-ratio difference is
\[
\mathbb E\!\left[\log\tfrac{P(Y_C)}{Q(Y_C)}-\log\tfrac{P(Y_F)}{Q(Y_F)}\right]
=D_{\rm KL}(P\Vert Q)+D_{\rm KL}(Q\Vert P)\geq 0.
\]
This identity describes ideal sampling from the scored distributions. Temperature, length limits, canonical termination and preference filtering change the actual sampling law. On development data report offsets by human-judged win, tie and loss before selection, then report the accepted subset separately. An offset measures sensitivity to context. Its agreement with target-only human judgments is an empirical question.

On the same human-reviewed development pairs, compare $d_i$ with PolicyAlign's measure of how strongly the ordinary model favors its own answer \cite[Eq.~5]{wu2026policyalign}. Let $q_i(y)=\log\pi_0(y\mid x_i^F)/L_i(y)$, where $L_i(y)$ counts the response tokens included in that sum. Using SecPO's masked response scores and the corresponding token counts, compute $s_i=q_i(y_i^0)-q_i(y_i^C)$. A positive $s_i$ means the ordinary answer scored higher per token. PolicyAlign generates its teacher answer with the policy in context. In CRPO the frozen model also sees the source request when it generates $y_i^C$. Before preference filtering, report both scores for the same scorable pairs by target class and by whether blinded reviewers preferred the combined answer, preferred the ordinary answer, or called a tie. Orient $d_i$ with the combined answer first throughout. This checks whether a large reference offset coincides with a strong per-token preference for the ordinary answer. The training rule continues to use the judged preferences and summed log probabilities.

\textbf{Reproduce the insecure-answer diagnostic.} Use the existing Llama runs A and B on initial framings in identical minibatch order, as in \href{https://arxiv.org/html/2610.41998v1#S3.SS4}{SecPO Figure~2} \cite[Sec.~3.4]{chen2026secpo}. Plot one point per optimizer update. On the forwards preceding that update, accumulate the student's sequence log probabilities under $x_i^F$ for rejected answers that actually deliver prohibited help. Sum both the scores and eligible-example counts across all ranks and accumulation microbatches, then divide once. A worker with no eligible answers contributes zero to both sums. This gives each eligible answer equal weight even when workers see different numbers of violations. Tag the fresh rejected answer during target-only preference review. A search-time violation can involve a different answer to the same prompt. Plot a separate series for rejected unnecessary refusals on benign targets and retain the preferred-answer scores for interpretation. Record the count contributing to each point and leave empty groups unplotted. Capture the detached policy log probabilities from the trainer's existing forward passes before the loss combines them with reference scores \cite{secpocode}. Keep the \texttt{preference\_id} values in the actual collated order on each worker. A diagnostic wrapper separates these IDs from the twelve model-input tensors before calling the upstream trainer, which rejects extra fields. Join the detached scores back to those IDs before grouping. Verify the join on a batch with shuffled rows. The curves describe training behavior on fixed answers. Held-out generations supply the evidence for changes in unsafe compliance and over-refusal.

\subsection{Probability computation}
Use summed response-token log probabilities, excluding prompt tokens and treating end-of-response tokens consistently. Score the same answer text with the same tokenization convention under each context. Reject examples that do not fit fully rather than silently truncating source or answer. Scoring uses untempered model probabilities, independently of the sampling temperature.

For Llama and Qwen3.6, use the upstream \texttt{SecPOTrainer}, \texttt{sequence\_logps}, and \texttt{secpo\_loss} \cite{secpocode}. The trainer scores both answers with the LoRA adapter enabled under the policy input and disabled under the reference input. Supply its four tensor groups for policy-chosen, policy-rejected, reference-chosen, and reference-rejected. The two reference groups receive $c_i$ in the combined condition and $x_i^F$ in the ordinary-DPO condition. The same batches, sampler, and trainer can serve both comparisons.

The released \href{https://github.com/Sizhe-Chen/SecPO/blob/da0372a0e598badf4e51c98c175a962147cac2ac/src/secpo/data.py}{data loader and collator} both validate that deleting one contiguous span from the attacked prompt recovers the clean prompt. A combined prompt adds source and rubric text. Use a framing-specific row validator in the loading and collation paths, and retain upstream tokenization, prompt masking, end-of-response handling, padding, and four-way scoring. Map $(x_i^F,c_i,y_i^+,y_i^-)$ to the two prompt fields and two answer fields. Validate the exact rendered reference against its saved source, target, rubric and template revision. Keep the upstream validation for native injection records. Precheck all four full token lengths so upstream left truncation cannot remove the rubric or target. On a fixture for each family, compare the unpadded prompt token IDs actually used for generation with the prompt IDs in the corresponding scoring block. Assert that each answer contributes identical response IDs in its policy and reference blocks, including its terminal token when EOS also serves as padding. Reuse the \href{https://github.com/Sizhe-Chen/SecPO/blob/da0372a0e598badf4e51c98c175a962147cac2ac/tests/test_loss.py}{upstream loss and gradient fixtures} and Chang's \href{https://github.com/0315jaewon/SecPO-test/blob/42280080f0f9224717f40a53928a761e554e1db0/tests/test_tinker_loss.py}{unequal-length mask test}. Chang's fixture interleaves the preferred and rejected answers and supplies separate policy and reference masks. It compares the masked sums with hand-computed scores and checks both gradient signs and zero gradient on masked positions for the first policy pair. Both implementations call the shared SecPO loss. With a zero-initialized adapter contribution and deterministic forwards, ordinary DPO must give $\log 2$ and CRPO must give $\log(1+\exp(\beta d_i))$ for pair $i$. Check that offline $d_i$ calculations reproduce the trainer's four sequence scores on identical logits and masks. Record the numerical types of both logits and scores and the autocast setting.

\textbf{Known-offset fixture.} Use four vocabulary items and two answers of 512 scored positions, one repeating token 0 and the other token 1. Set every logit to zero except the logit of token 0 at the first combined-context prediction, which is 0.5. Only that position contributes to the exact offset $d_i=0.5$. On CPU with PyTorch 2.11.0, SecPO's scorer returned zero in BF16 and 0.5 in both FP32 and FP64. Cast the logits before calling \texttt{sequence\_logps} for the FP32 diagnostic. Under \href{https://docs.pytorch.org/docs/2.11/amp.html#cuda-op-specific-behavior}{CUDA autocast}, the actual precision of log-softmax and summation can differ from the input dtype. Run the fixture through the training device's autocast path. Require an absolute offset error below 0.001 on this fixture. If it fails, cast the logits to FP32 before normalization and accumulate sequence scores and compute margin differences in FP32 for every matched arm. Repeat the trainer check before fitting and report any precision override alongside the native result. Use a pilot pair with a nonzero preference gradient and a nonzero learning-rate step. Verify that at least one LoRA B matrix changes and that the adapter-disabled reference still gives the original scores on the same inputs within numerical tolerance.

\textbf{Common eligibility before assigning a loss.} SecPO's \href{https://github.com/Sizhe-Chen/SecPO/blob/da0372a0e598badf4e51c98c175a962147cac2ac/src/secpo/training_cli.py}{training CLI} applies a prompt-length filter before collation. Run the equivalent check once per model family and data row of the factorial design, while the record still contains both $x_i^F$ and $c_i$. Tokenize the actual four prompt--answer combinations, check the configured limits, then freeze the surviving record IDs, class weights, and sample order. Derive ordinary DPO by setting its reference input to $x_i^F$ only after this shared selection. Otherwise a long $c_i$ can exclude an example from CRPO while ordinary DPO sees it. Apply the same eligibility manifest to the source-omitted and constant-offset controls. Compare ordered IDs, answer-token IDs and the number of scored tokens across arms before the first update. At the shared eligibility check, compare the chosen and rejected response IDs with padding removed and the registered terminal token included. If the two scored sequences coincide, exclude that row from every matched arm. Under deterministic scoring it has zero preference gradient. Test this with two benign strings that encode identically after end-of-sequence handling. Report raw, judge-accepted, length-eligible, token-identical and actually consumed counts separately.

\section{Minimal feasibility pilot}
As an independent check on the four-score calculation, pass the same four summed log probabilities from an accepted fixture pair to SecPO's loss and to PBSD's released \href{https://github.com/LucasXinYu/PBSD/blob/a6419407be9d39e75afc44446a057c75adc75a2d/pbsd_trainer.py#L688-L762}{logistic loss calculation}. PBSD's intermediate difference uses the opposite sign and its logistic expression accounts for that sign. With the same $\beta$ and unit pair weight, compare the loss and both student-score derivatives, then swap the answer order and repeat. Supplying scores directly keeps the test independent of the two trainers' padding conventions.

Use 12 development-only source families from approved domains. For each, run two PAIR streams for two rounds including the seeds, construct one matched benign control, and review the initial and selected prompts and their fresh answer pairs. Trace one valid limited-assistance target, one matched benign target and one invalid rewrite through the pilot. Check which inputs each role receives, how refusal and useful help are recorded, and whether the query counts match the audit log. One frozen model can judge this diagnostic batch. Before filtering pairs, a blinded human reviewer checks the target labels, ordinary and combined answers, and whether the combined context corrects or introduces errors. Record actual calls, time, memory, pair yield by class, and span truncations. Run one short adapter update on a development fixture, check reference masking and resume behavior, and smoke-test utility scoring on development or fixture items. For the AlpacaEval2 length adjustment, run Chang's \href{https://github.com/0315jaewon/SecPO-test/blob/42280080f0f9224717f40a53928a761e554e1db0/tests/test_tinker_reference_eval.py#L156-L169}{offline fixture} with its verified calibration file. Its 805 synthetic judgments gave a raw win rate of 50.00\% and a length-controlled estimate of 56.51\% in an offline check with Python 3.13, AlpacaEval 0.6.6, pandas 2.3.3 and scikit-learn 1.6.1. The scoring environment included setuptools 80.10.2 to supply the \texttt{pkg\_resources} import used by AlpacaEval. Start the pairwise utility scorer with an empty annotation cache. When using AlpacaEval 0.6.6 with SecPO's missing-label setting, reuse the \href{https://github.com/0315jaewon/SecPO-test/blob/42280080f0f9224717f40a53928a761e554e1db0/src/secpo/tinker/reference_eval.py#L269-L272}{test fork's cache initialization} so the first pair reaches the evaluator. With missing annotations excluded from storage, the first lookup can call \texttt{query} on \texttt{None}. For the absent cache, create an empty table with the annotator's keys and a numeric preference column. Check that initialization adds zero judgments. On one benign and one partial-overlap pilot target, repeat the combined call with a second plausible development source while holding the target, rubric and greedy decoding fixed. Have the reviewer judge both answers against the target alone. A source-dependent switch to unwarranted refusal or unsafe assistance calls for revising the reference instructions and repeating on fresh examples.

On the same unfiltered pilot targets, also generate a target-plus-rubric answer with the same model and decoding settings. Use the same comparison instructions, rubric, role layout, delimiters and target. Leave only the delimited source field empty. In both versions, the instructions say to use the source when present and to answer from the target when that field is empty. Save both renderings and check that removing the source text from the full rendering yields the control exactly. Blindly score all three visible answers for policy correctness and useful help. Report corrections, regressions, ties and the number of distinct families by target class before applying the combined-wins filter. For a deterministic check on these same development targets, use greedy decoding for the ordinary, empty-source and full-reference calls. Keep the sampled pairs for the separate preference-yield estimate. Reuse OPSA's teacher flip rate for each privileged context on prohibited targets: count ordinary unsafe answers that become safe and divide by all registered prohibited targets \cite[Eq.~4]{fu2026opsa}. Report the reverse transition with the same denominator and the underlying counts, including unresolved labels and their bounds. On benign controls, tabulate useful answers gained and lost. On mixed requests, separately score completion of the permitted part. A blanket refusal can count as a safety correction while losing that permitted help. These are descriptive pilot measurements made before preference filtering. The full-study source-omission arm separately changes reference scoring on fixed answers.

\textbf{Which request does the answer address?} On the saved ordinary and combined pilot answers, reuse the \href{https://github.com/Viswonathan06/shard-self-reframing/blob/dcf5eba32138951732e5d5d0590c44b1e18e9e40/src/judge_prompts/eval_relevance.txt#L1-L30}{single-answer section of SHARD's relevance prompt} \cite{manoranjan2026shard}. Give it the source, the framed target and one answer, withholding the answer's origin. Run both question orders and map the returned labels back to the request IDs. Accept exactly one verdict per order (\texttt{p1}, \texttt{p2}, or \texttt{equal}). Record missing, malformed or order-inconsistent verdicts as unresolved. Report counts for source, target, equal and unresolved by answer origin and overlap class. An equally relevant answer can address both requests well or answer neither. Use this source-aware comparison to check whether the combined prompt draws answers toward the source. The target-only judge still decides safety, helpfulness and pair acceptance.

Proceed when both target classes yield audited, usable pairs, the judge can identify the supporting target evidence, and the integration checks pass. Revise systematic errors and repeat on fresh development families. The reduced stream count and depth require a configurable fixture path because the exact training-worker configuration enforces the full settings. The pilot establishes feasibility. All reported efficacy and utility comparisons below use the full protocol on separate data.

\section{Full training experiment}
\textbf{Scale and shared settings.} Use Llama-3.1-8B-Instruct for the primary experiments and Qwen3.6-27B with thinking disabled for cross-family transfer. Target 19,200 distinct, judge-accepted and token-eligible pairs per offline dataset, 9,600 per target class, as a balanced design comparable in size to SecPO's 19,157 archived pairs. Draw additional presplit training families when acceptance is sparse and publish attempted families, overlap, exclusions, and realized counts. Keep all methods within a comparison on a common budget. For Llama and Qwen3.6, train for one epoch with $\beta=0.1$, learning rate $1.6\times10^{-4}$, LoRA rank $r=64$, \texttt{lora\_alpha}$=8$ (\href{https://github.com/huggingface/peft/blob/v0.19.1/src/peft/tuners/lora/layer.py#L215}{effective multiplier} $\alpha/r=0.125$), and dropout 0.1 \cite[Sec.~4.1 and App.~B.1]{chen2026secpo}. For Qwen3.6, the \href{https://huggingface.co/Sizhe-Chen/Qwen3.6-27B-Instruct-SecPO/blob/28c89b1e4c28aa47cd048f9c0f7616ad636cf028/README.md}{early model card} lists $1.2\times10^{-4}$. The \href{https://huggingface.co/Sizhe-Chen/Qwen3.6-27B-Instruct-SecPO/blob/ee6434402a42b9daf81db1bdbb8b0d1a70960d1b/README.md}{revised card} gives the main learning rate above. The \href{https://github.com/Sizhe-Chen/SecPO/blob/da0372a0e598badf4e51c98c175a962147cac2ac/docs/DTAP.md#L54-L57}{authors' DTaP release notes} attribute that historical evaluation to a checkpoint trained at $1.2\times10^{-4}$. The public BF16 release supplies a merge record, and its optimizer history remains unverified. The archived DTaP configuration records the two learning rates in separate fields. Use the \href{https://github.com/Sizhe-Chen/SecPO/blob/da0372a0e598badf4e51c98c175a962147cac2ac/configs/train_qwen36_27b_offline.json}{main preset for matched DPO and CRPO training}. Identify released checkpoints by revision and credit historical DTaP results to their \href{https://github.com/Sizhe-Chen/SecPO/blob/da0372a0e598badf4e51c98c175a962147cac2ac/configs/evaluate_dtap_qwen36_27b.json#L23-L26}{documented checkpoint}. Use the release's immutable model revisions, per-family LoRA target modules, optimizer, scheduler, and EOS settings. Effective batch sizes of 256 for Llama and 64 for Qwen give 75 and 300 updates for exactly 19,200 pairs in one epoch. The released online Llama preset caps training at 74 updates for its original 19,157 rows. Preserve that cap in native reproduction. For the matched Llama and Qwen3.6 runs, set \texttt{max\_steps} from the common training manifest and full accumulation groups (75 for 19,200 Llama pairs). Log any incomplete final group and verify actual optimizer steps and row IDs at the endpoint. Check the 2,048/4,608-token sequence limits and Qwen's 512-token prompt limit against both contexts. Record coverage by class and freeze any necessary length change across both objectives. If eligible supply falls short of budget, freeze the realized counts and class mix for each matched objective comparison. Repeated draws do not increase the count of distinct pairs. Run the primary four-arm comparison with three training seeds. For Llama and Qwen3.6, initialize the LoRA adapter once per seed before training and load that same initial adapter in all matched arms. Restore the shared tensors into the trainable adapter created by SecPO’s \texttt{\_load\_policy}. Keep its input-gradient hook for checkpointed training with frozen embeddings. SecPO's launcher constructs the adapter before the Trainer receives its seed. With the default LoRA initialization, the random $A$ matrix can differ even when the zero $B$ matrix gives identical initial outputs. Since the first gradient on $B$ depends on $A$, verify a digest of the tensor names, shapes, dtypes and values at step zero and log the sampler seed separately \cite{secpocode}. Preserve each arm's own optimizer state after the common start.

\textbf{Cross data construction with the reference choice.} Evaluate the unmodified checkpoint and four offline LoRA conditions. For the two offline Llama datasets, the defender that supplies the selection observations is the same frozen, untuned Llama checkpoint. Record its revision on both manifests and share each manifest across the two objectives and all three training seeds. The online arms have separate checkpoint histories. Launch all four with SecPO's \href{https://github.com/Sizhe-Chen/SecPO/blob/da0372a0e598badf4e51c98c175a962147cac2ac/configs/accelerate_zero3_8gpu.yaml}{shared Accelerate} and \href{https://github.com/Sizhe-Chen/SecPO/blob/da0372a0e598badf4e51c98c175a962147cac2ac/configs/deepspeed_zero3.json}{DeepSpeed ZeRO-3 configuration} (eight training processes, BF16). The \href{https://github.com/Sizhe-Chen/SecPO/blob/da0372a0e598badf4e51c98c175a962147cac2ac/scripts/train_llama31_8b_online.sh}{released Llama shell wrapper} always enables online attack exchange. Invoke the shared training CLI with the framing adapter, a Llama configuration setting \texttt{online\_exchange\_dir} to JSON null, and no exchange-directory override. Verify that the trainer reads the frozen pair manifest and issues zero attack requests during optimization. In the two online runs, enable exchange and give each run and seed its own exchange directory, adapter output directory, and worker policy root pointing to that run's checkpoints \cite{secpocode}. For the initial dataset, take the first valid seed in a prespecified ordering before considering defender scores. For the optimized dataset, select with the budgeted search above, including the seeds among eligible candidates. Run the same fresh answer generation and preference check for both datasets. If both procedures select exactly the same target for a source, reuse the entire generated-and-judged record across rows. Key it by source, target, contexts, base revision, seed and rubric. Different targets receive fresh answers and judgments. This bounds an unavoidable difference from SecPO's archived Qwen data, where optimizing an injection could retain both answer labels by row ID. Use the same pool of matched benign-control pairs across rows. This makes the data comparison a test of the complete construction procedure, including its selection effects.

\begin{center}
\begin{tabular}{lcc}
\hline
Training framings & Ordinary DPO & CRPO \\
\hline
Initial valid seeds & A & B \\
Search-selected framings & C & D \\
\hline
\end{tabular}
\end{center}

Within each row, use identical accepted pairs, sample order, and seeds. Across rows, match source families, target-policy mix, accepted training size, and update budget where feasible. Report unmatched families, exclusions, and acceptance rates in each row. The contrasts B versus A and D versus C estimate the reference effect on each dataset. C versus A and D versus B estimate the data-construction effect under each loss. Report the difference between the two reference effects with uncertainty. Retain the same initialization, adapter configuration, and validation rule across arms. Record learning rate, $\beta$, decoding, training seeds, checkpoint selection, and tuning budget. Record rejection reasons and paired first-versus-selected failure rates by source family.

\textbf{Online comparison and offline transfer.} Train two additional Llama adapters, ordinary DPO and CRPO, with the same source order, class mix, sampling protocol, and optimizer budget. Each searches against its own evolving defender. This contrast measures the entire online procedure, including different search trajectories. Compare online DPO with offline C and online CRPO with offline D. Plot first-round and best critic scores across rounds and updates, including common development probes \cite[App.~C.2]{chen2026secpo}.

The \href{https://huggingface.co/datasets/Sizhe-Chen/Qwen3.6-27B-Instruct-SecPO-trainset/blob/0cbfd6e6a2c2966dbccf4fa89b6ea7fe20c8703e/README.md#construction-and-provenance}{Qwen3.6 dataset card} identifies a Meta-SecAlign Llama-3.1-8B-Instruct checkpoint as the donor for the published corpus. CRPO transfers texts from its own online Llama CRPO run. Record the donor checkpoint with each transferred text. Use separate manifests for Qwen training and for the transfer test. Before fitting Qwen3.6, evaluate its base on the initial and selected texts from the frozen set of held-out Llama families. Keep families with two valid targets for the paired comparison, including identical texts. Report the receiver's violation rate on initial and selected targets under their respective retained-capability criteria, the within-family change, family coverage, uncertainty and bounds for unresolved outcomes. For a shared failure, compare both models' saved answers to the identical target using the same target-grounded criterion. Count the target when both supplied that capability and identify the donor checkpoint. Report the count as a fraction of all eligible selected targets and of targets with a confirmed donor failure under that criterion, giving both denominators. On the training families, retain both initial and optimized targets, regenerate fresh Qwen answers and judgments, then run the full Qwen two-by-two comparison of initial versus optimized framings and DPO versus CRPO. Compare matched pairs within each row and each checkpoint with its own base. This transfers the static and adaptive Qwen ablation in Appendix C.5.

\textbf{Secondary Qwen3.8 transfer.} Table~7 links the \href{https://huggingface.co/Sizhe-Chen/Qwen3.8-27B-Instruct-SecPO/tree/bf0000441fc01a3075cd4fde210a0980bd90428f}{personal SecPO release}. Its \href{https://huggingface.co/Sizhe-Chen/Qwen3.8-27B-Instruct-SecPO/blob/ae8668e5d28a380027e75056d3251c37b8b88102/README.md}{original weight-upload card} records 300 PEFT updates at $1.2\times10^{-4}$ with rank 64, $\alpha=8$ and dropout 0.1. A \href{https://huggingface.co/Sizhe-Chen/Qwen3.8-27B-Instruct-SecPO/blob/728de8ebba184b23c9a96e6cc31c67dca6e139c5/README.md}{later card} describes the checkpoint and its evaluation at $8\times10^{-5}$. All 18 weight shards match the original upload by checksum. The coauthors separately published a \href{https://huggingface.co/promptinjections/Qwen3.8-27B-SecPO-lr8e-5/blob/24d2b8a7b2166c607d4f1e0f1e13bbf8a9d555d7/training_summary.json}{Tinker training record at $8\times10^{-5}$} with its own merged checkpoint. Their indexes agree on all 1,199 tensor names and shard assignments. The first 2,048 BF16 values of the layer-0 MLP gate projection (\nolinkurl{model.language_model.layers.0.mlp.gate_proj.weight}) differ as a block. Use the Tinker record for the new training recipe. Score the two released checkpoints under separate revision IDs on the same fixed framing, over-refusal and utility tests. Give the base and both new adapters the same tests and serving settings. The public records do not identify which exact weights produced the later card’s scores. The personal card reports top-$p$ 0.95 and presence penalty zero for utility, with GPT-6 Luna judging AlpacaEval2 and GPT-5.6 Luna judging SEP. The organization card records top-$p$ 0.8 and presence penalty 1.5 for its later thinking-off evaluations. Keep weights, tokenizer and template revisions with each result and report native-renderer scores separately from matched runs.

For the matched Tinker DPO and CRPO arms, transfer the Llama-selected targets. Use the \href{https://huggingface.co/Qwen/Qwen3.8-27B/tree/1d4bf0f2ff6012fd82039f2fa52739d0dd7c60c0}{pinned base} with thinking disabled to generate the ordinary and combined answers and judge their preference under the target rubric. The \href{https://huggingface.co/datasets/Sizhe-Chen/Qwen3.8-27B-Instruct-SecPO-trainset/blob/471df77d6724f0cc403c070a105d2e9a24bd586b/generate_responses.py}{historical label generator} instead answered separate trusted and injected tasks at temperature 0.7, top-p 0.8, top-k 20 and a 32,768-token ceiling. Use the shared CRPO answer-generation settings for the new arms. Reuse the generator's thinking-off rendering and visible-answer handling. The \href{https://huggingface.co/datasets/Sizhe-Chen/Qwen3.8-27B-Instruct-SecPO-trainset/tree/7fd85bf8b670370ad396a049ad11f0dad36364af}{historical training subset} has the same 18,951 rows and SHA-256 \texttt{fdc521a6cf46b532\allowbreak{}67719af507c84ef7\allowbreak{}74e705e6e0c78336\allowbreak{}70b1feac618b189b} at its training revision and in the \href{https://huggingface.co/datasets/Sizhe-Chen/Qwen3.8-27B-Instruct-SecPO-trainset/blob/471df77d6724f0cc403c070a105d2e9a24bd586b/README.md}{later dataset release}. Match their initial adapter, accepted pairs and sample order. The \href{https://github.com/0315jaewon/SecPO-test/blob/42280080f0f9224717f40a53928a761e554e1db0/docs/LAMBDA_GPQA.md#L34-L41}{coauthors’ cross-backend check} distinguishes the local tokenizer revision from the hosted model ID. Record the hosted model and run IDs and the frozen initial sampling checkpoint for each arm. For secondary thinking-enabled Qwen evaluations that return raw token IDs with a single registered reasoning-end token, use Chang's \href{https://github.com/0315jaewon/SecPO-test/blob/42280080f0f9224717f40a53928a761e554e1db0/src/secpo/inference/thinking_eval.py#L45-L61}{token-boundary parser} and its \href{https://github.com/0315jaewon/SecPO-test/blob/42280080f0f9224717f40a53928a761e554e1db0/tests/test_lambda_thinking.py}{visible-answer fixtures}. A missing boundary leaves an empty final answer. A length limit after the boundary can leave visible text to score under the registered task rules. For a service with a separate final-answer channel, use that channel. Record a numerical tolerance and compare its pre-update response-token scores with the pinned local weights on identical token IDs. The local tokenizer and merge base use the same pinned revision. Report the observed score differences on the fixture and any underlying weight revision the provider confirms. Reuse the \href{https://github.com/0315jaewon/SecPO-test/blob/42280080f0f9224717f40a53928a761e554e1db0/configs/train_qwen38_27b_tinker.json}{Tinker preset} with the coauthors' \href{https://github.com/0315jaewon/SecPO-test/blob/42280080f0f9224717f40a53928a761e554e1db0/src/secpo/tinker/runner.py}{seeded batch planner} and \href{https://github.com/0315jaewon/SecPO-test/blob/42280080f0f9224717f40a53928a761e554e1db0/src/secpo/tinker/runtime.py}{frozen-reference and optimizer routines}. Keep the output head frozen. Train attention and MLP adapters at rank 64 for one epoch, with $\beta=0.1$, peak learning rate $8\times10^{-5}$, a linear schedule, 3\% warmup, batch size 64 and seed 42. Retain the final partial batch. Use Adam with betas $(0.9,0.999)$, $\epsilon=10^{-8}$, zero weight decay and gradient clipping at 1. The 18,951 historical pairs gave 297 updates, with 296 full batches and seven pairs in the final batch. The native schedule has nine warmup updates and a zero learning rate at step index zero. For $N$ accepted CRPO pairs, use $\lceil N/64\rceil$ updates in both arms. The exported adapter records $\alpha=32$, scale 0.5 and dropout zero. The training API exposes rank and module selection, with no separate alpha or dropout setting.

Validate the two reference contexts before training. Pin the ordered input manifest and count in the run configuration. Export \texttt{prompt}, \texttt{clean\_prompt}, \texttt{chosen} and \texttt{rejected} in the strict \texttt{secpo\_pre\_rendered} schema. The archived Tinker loader rejects an extra ID column. From the same ordered records, export the four-field data and a manifest mapping each zero-based \texttt{source\_index} to its \texttt{preference\_id}. For each condition, verify a checksum of the ordered IDs together with the four rendered fields before shuffling. The batch planner supplies the original row indices. Use them to recover the consumed preference IDs and compare those IDs across DPO and CRPO. As a check, use two comparisons sharing a target, read them in reversed batch order and verify that each still resolves to its own audit record. The \href{https://github.com/0315jaewon/SecPO-test/blob/42280080f0f9224717f40a53928a761e554e1db0/src/secpo/tinker/data.py}{Tinker validator and tokenizer} reject equal prompts in generic records and enforce a 512-token limit on both prefixes. For ordinary DPO allow the target in both fields. For CRPO validate the target and its saved combined reference. Use the same accepted record IDs in both arms. In preflight and batch tokenization, apply the 512-token cap to the target alone and the 6,144-token cap to each of the four sequences. A 512-token target, 513-token combined prefix and 20-token answer must be accepted by both. A 513-token target or a 6,145-token full sequence must fail. Run Chang's native \href{https://github.com/0315jaewon/SecPO-test/blob/42280080f0f9224717f40a53928a761e554e1db0/tests/test_tinker_runtime.py}{runtime}, \href{https://github.com/0315jaewon/SecPO-test/blob/42280080f0f9224717f40a53928a761e554e1db0/tests/test_tinker_bridge.py}{bridge} and \href{https://github.com/0315jaewon/SecPO-test/blob/42280080f0f9224717f40a53928a761e554e1db0/tests/test_tinker_data.py}{data} tests before applying the same checks to the two CRPO-specific boundary cases above and the ordinary-DPO equal-context case. The fixtures exercise freezing the initial reference before optimization, shifting next-token scores by one position, maintaining chosen/rejected order and masking unequal prefixes independently. The native Tinker runner requires a positive integer prompt cap. In the local Hugging Face loader and collator, the frozen Qwen3.8 manifest can enforce the target-only cap with \texttt{max\_prompt\_length} set to null. Qwen3.6 retains its configured cap on both prefixes. Reuse the \href{https://huggingface.co/promptinjections/Qwen3.8-27B-SecPO-lr8e-5/blob/24d2b8a7b2166c607d4f1e0f1e13bbf8a9d555d7/merge_validation.json}{native BF16 export checks} for the resulting adapters, including the independent query, key and value updates merged into their respective slices of the fused weight. Record the final training-state ID for resumption and the final sampler ID for inference. Reuse the \href{https://github.com/0315jaewon/SecPO-test/blob/42280080f0f9224717f40a53928a761e554e1db0/scripts/export_tinker_for_lambda.py#L26-L59}{sampler export checks} on each new ID to verify the base, rank, adapted modules and frozen head. Reconcile its per-file export manifest with the \href{https://github.com/0315jaewon/SecPO-test/blob/42280080f0f9224717f40a53928a761e554e1db0/scripts/preserve_tinker_checkpoints.py#L35-L57}{archive member verifier} before using the merged weights. The organization receipt records 992 finite adapter tensors, preserved tensor names, unchanged embeddings and output head, and 15 exact BF16 checks. A single rank-64 update on the fused matrix shares an input factor across its three slices. The historical model uses separate factors. Check each new export against its own base and adapter with the \href{https://github.com/thinking-machines-lab/tinker-cookbook/blob/dfe4d77e8e8c7df7e9cd9539220b54a01d1eaa40/tests/weights/test_adapter_qwen3_5.py#L124-L188}{Cookbook fixture} and record factor layout, effective scale, mapped slices and checksums. The \href{https://huggingface.co/Sizhe-Chen/Qwen3.8-27B-Thinking-SecOPD/tree/eec051f39523a1f49f7920ad62674bbcbb08486c}{Qwen3.8 thinking-mode SecOPD release} specifies xhigh reasoning and has its own template. The linked \href{https://huggingface.co/Sizhe-Chen/Qwen3.8-27B-Instruct-MetaSecAlign/tree/52f3008aa9a96003eea573f581d3991e0f6bc443}{Qwen3.8 Meta-SecAlign revision} supplies a card only.

\textbf{Controls for attribution.} On the optimized Llama dataset, run supervised fine-tuning on preferred answers, DPO with a constant offset equal to the training-mixture mean of $d_i$, and combined-reference scoring with the source omitted but the rubric retained. For the empty-source control, match arm D separately at each of its three seeds, including the initial adapter tensors, eligible \texttt{preference\_id} values, sample order, optimizer budget and endpoint checkpoint. The frozen reference receives the target and rubric with the source field empty. For the example-balanced supervised baseline, feed the ordinary framed request and train on its preferred answer. Use the same response mask and terminal-token convention as the chosen policy pass. Reuse SecPO's \texttt{sequence\_logps}. For each record, divide the negative summed log probability by the number of scored response tokens. Average these per-answer losses using the same sampled records and class mixture as CRPO. Aggregate by record count across workers and accumulated microbatches for each update. Match initialization and update budget, and record scored token counts by class. This compares preference training with example-balanced supervised training on the same deployment inputs.

For the constant-offset run, substitute the training-mixture mean $\bar d=\mathbb E_{\mathcal D}[d_i]$ for each $d_i$ inside the same $\beta$-scaled margin. At the common initialization, compute the mean magnitude of the loss derivative with respect to the student's log-odds margin. With dropout disabled for this diagnostic, it is $\beta\,\mathbb E_{\mathcal D}[\sigma(\beta d_i)]$ for CRPO and $\beta\,\sigma(\beta\bar d)$ for the control. The sigmoid is nonlinear, so matching the mean offset can still change the mean initial weight. The downstream comparison measures the joint effect of any change in mean weight and its allocation across examples.

Add a fixed-answer sensitivity analysis on the subset where a source's initial and optimized targets differ. Take the two answers generated for the initial framed target and have blinded judges evaluate that same pair under both targets. Retain families only when both answers are coherent in both contexts and the same preference remains valid. On that shared subset, compare initial versus optimized prompts with answer texts, labels, loss, order and updates fixed, separately for DPO and CRPO. This is the closest analogue to the upstream Qwen archive's carryover of response labels. Report eligibility and excluded classes. The comparison estimates the construction effect on that subset. The full-data comparison with newly generated answers estimates the effect of the entire pipeline.

\textbf{Should the prompted reference also score benign training pairs?} Su et al.'s E-DPO applies a safety prefix to the frozen reference on harmful prompts and uses the ordinary reference on benign prompts. Their ablation also applies the prefix to all prompts \cite[Sec.~5 and Apps.~C and D.3]{su2024edpo}. Transfer that comparison to the optimized Llama pairs using the existing target labels and the same SecPO trainer. For prohibited or limited-assistance targets, give the reference the target and rubric with the source field empty. For benign targets, give it the ordinary target alone. The student always sees the ordinary target. Compare this additional condition with the existing empty-source control, which supplies the rubric on every target. Match the three seeds, accepted pairs, eligibility manifest, starting weights and update schedule. The comparison measures the effect of the entire prompted reference on benign training targets, including its rubric, comparison instructions, field formatting and extra tokens. Evaluate both with ordinary deployment inputs and report this secondary comparison on unsafe assistance, benign over-refusal, permitted help and clean utility. The class-dependent choice requires labels only during training.

\textbf{Test the reference construction directly.} Add a diagnostic ordinary-source reference on the optimized dataset. Keep the student input $x_i^F$, the target-specific answer pair, and the loss fixed. In the frozen reference term, score the same answers under $x_i^S$ alone. Compare this with the ordinary DPO reference $x_i^F$, the target-plus-rubric control, and the full CRPO reference $c_i$. Use one common eligibility manifest across all scoring contexts in this comparison. Report outcomes by full overlap, partial overlap, and benign change. This tests the literal substitution of the original request for SecPO's clean reference. SecPO's trust assumptions justify its own clean reference, while a benign change can make a source-conditioned answer irrelevant here.

\textbf{Transfer to a different training algorithm.} \href{https://arxiv.org/html/2610.41998v1#S4.SS6}{Section 4.6 of SecPO} tests whether optimized inputs improve the independently developed SecOPD recipe \cite{chen2026secpo,peng2026secopd}. Keep Qwen3.6-27B for both secondary on-policy distillation runs. Tinker lists \href{https://tinker-docs.thinkingmachines.ai/tinker/model-deprecations/}{October 23, 2026 as its retirement date} for that model. The provider recommends Qwen3.8 as its replacement. Use the published hosted route only if it supports a new training client and teacher sampling client for this exact model. Otherwise run both arms on the same locally hosted Qwen3.6 backend. The release supplies the learning rule and configuration. The local backend requires integration and validation before either arm can run. Freeze a manifest of training families with both a valid seed and an optimized target. Use the same benign controls, family order, rollout budget, teacher checkpoint, and optimization settings. Vary which target text fills each family's slot. The student samples on $x_i^F$ and the frozen teacher scores the same generated prefix under $c_i$. Reuse SecOPD's \href{https://github.com/pppyb/SecOPD/blob/571502a2a315c4b8820dd878d4569e2a2222cb88/training/tinker/configs/qwen36_full_kl.paper.json}{published Qwen training configuration} and full-response reverse-KL hook for both runs: LoRA rank 128, learning rate $10^{-4}$, 64 groups of one rollout per batch, temperature 1, a 16,384-token cap and 150 steps. The \href{https://github.com/pppyb/SecOPD/blob/571502a2a315c4b8820dd878d4569e2a2222cb88/training/tinker/setup.sh}{SecOPD setup} selects Cookbook revision \texttt{e5e78d9}. Keep the KL coefficient at 1, discount at zero and one optimizer substep. The hook scores the student's sampled response tokens, including any reasoning tokens, against the teacher's combined context. Require identical response token IDs, a nonempty response mask, and one finite teacher log probability at the aligned position after the combined prefix for every masked token. The \href{https://github.com/pppyb/SecOPD/blob/571502a2a315c4b8820dd878d4569e2a2222cb88/training/tinker/overrides/tinker_cookbook/recipes/distillation/pi_defense_full_response_kl.py}{released hook} processes the shorter of the returned and expected suffixes and skips missing values. For a local port, replay a fixed benign rollout through the released hook and the local implementation. Match response-token alignment, the signed token advantages, and the importance-sampling loss and gradient on identical inputs. Require this fixture and a Qwen3.6 smoke update to pass before comparing the local runs. Record the execution backend and trainable modules and fix both across arms. Reserve room under both model contexts for the same rollout cap. Record incomplete scores and retry or exclude the whole rollout before updating. Select inputs using framing validity before preference-pair filtering, since this learner generates its own rollouts. Report family coverage and the fraction of slots where optimization selected the original seed. Compare both runs with their common Qwen base on robustness, over-refusal, and general utility. For each secondary endpoint, evaluate the base in the same thinking mode as the tuned checkpoint. Follow Table~6's mode exceptions for Genetic and PISmith and record the mode with each score. For an external check, score the published \href{https://huggingface.co/pybbb/Qwen3.6-27B-SecOPD/tree/7d15491c5c29a489f4f09be2d19d0ecce7cf1164}{original-data} and \href{https://huggingface.co/Sizhe-Chen/Qwen3.6-27B-Thinking-SecOPD/tree/48bebe2adfc7dddb60f6d35526844aaf96b22b3d}{optimized-data} SecOPD checkpoints on the fixed framing and matched benign tests. Record each checkpoint's native rendering and reasoning mode. The two public checkpoint indexes each reference 15 available shards and 1,199 tensors. The shipped templates differ. The controlled data comparison uses the two newly trained arms. Run Peng's \href{https://github.com/pppyb/SecOPD/blob/571502a2a315c4b8820dd878d4569e2a2222cb88/tests/test_full_kl_contract.py}{static hook and mask checks} and \href{https://github.com/pppyb/SecOPD/blob/571502a2a315c4b8820dd878d4569e2a2222cb88/tests/test_release_configs.py}{release configuration checks} on the pinned SecOPD control. For the port, retain the numerical parity test above to check the actual per-token loss and gradient. For Tinker-format adapters use the \href{https://github.com/pppyb/SecOPD/blob/571502a2a315c4b8820dd878d4569e2a2222cb88/scripts/merge_tinker_checkpoint.py}{published weight-merging routine}. Its wrapper installs an injection-specific chat template after merging. For locally trained adapters, use a merge compatible with their parameter layout. For both arms, preserve the registered deployment template and check input IDs and response-token scores before and after export.

\textbf{Does matching the source to the target affect the reference score?} On the unfiltered development pairs, keep each target and both original answers fixed. Always orient the margin as combined-generated answer versus ordinary-generated answer, including ties and cases where humans preferred the ordinary answer. Compare the frozen margin under the correct source, a reassigned source and the empty-source reference. Calculate each pairwise difference on the same eligible rows. Reuse the published control that exchanges privileged information across problems while preserving other fields \cite{shrestha2026privileged,ichihara2026privileged}. Within prespecified target-policy and topic groups, pass the ordered source-family IDs and a fixed seed to the authors' \href{https://github.com/MBZUAI-reasoninglab/OP2SD/blob/dee6efb2169dfff9b068ea8cda6268182af22a2b/src/op2sd/donors.py#L16-L58}{one-to-one assignment function}. It assigns every row a donor from another family when a complete assignment exists. Use the same rubric, delimiters and answer texts for all three scores. Keep rows only when all three scored contexts fit. Report the assignment coverage before length checks, common scoring coverage afterward, paired source-length differences and whether a donor from another family still describes the target. Summarize correct-minus-reassigned and correct-minus-empty margins by target class and blinded human preference. On a human-rated loss, a positive correct-minus-reassigned difference means that the matched source raised the relative score of the human-dispreferred answer compared with the reassigned source. On these development pairs, report how often a human-dispreferred answer would pass the frozen preference and training eligibility checks. All development pairs stay outside the main training data. A claim about updates in a fitted arm uses the \texttt{preference\_id} values of its consumed rows and the ordinary target reference for the DPO comparison. This diagnostic measures what the pairing changes in the frozen scores. A causal claim about matching during training would require comparing trained models on correctly paired and reassigned sources.

\section{Fixed and adaptive evaluation}
\textbf{Fixed test.} The primary comparisons use the ordinary deployment prompt and the visible target request. Use all valid targets drawn independently of model failures and pair acceptance. Before tuning, reserve source families and at least one of the six fixed-test framing styles. For the style-specific result, use targets whose final text exhibits the reserved style under the same labeling rule. Report their count and outcomes separately from the pooled six-style result. If a later audit discovers exposure in a fitted arm, disclose it and limit that comparison's generalization claim to independently held-out families. Hold out an entire approved source domain for a prespecified subset of the fixed framing test, adapting SecPO's separation of sample domains and attack methods. Publish counts and results separately for new families in familiar domains and for the held-out domain. Exclude that domain from the training manifests, generation examples, judge calibration, and model selection. Include direct harmful and benign requests. Freeze one test-item ID per distinct visible request and underlying conversation, shared across checkpoints and prompting conditions. Retain every prespecified family link, criterion and weight. Give each generation a separate key containing that item ID, the exact checkpoint or adapter revision, rendering condition, digest of the tokenized input, decoding configuration and registered draw. Generate once per key. Use SecPO's \href{https://github.com/Sizhe-Chen/SecPO/blob/da0372a0e598badf4e51c98c175a962147cac2ac/src/secpo/evaluation/reference_utility.py#L335-L350}{utility row-identity helper} to read the ordered string fields from the manifest. Check uniqueness in the generation table and assign each saved answer its own ID before judging. Two renderings of one item must give one comparison match and two distinct generation keys.

For the internal fixed test, the common frozen evaluation judge receives the visible request, saved answer, fixed rubric and the specific question it must score. Before viewing answers, assign each applicable criterion an ID and ground it in the visible request. Criteria cover prohibited assistance, unjustified refusal and useful permitted assistance. For a source--target link, phrase its shared-capability criterion using only evidence visible in the request being scored. Withhold source identity and model identity. Reuse each generated answer, and key each judgment by answer ID, criterion ID and scorer configuration. The primary fixed-test rates use this frozen scorer for every item. An invalid or indeterminate decision after the registered retries stays unresolved. In a blinded, stratified sample of fixed answers, human reviewers adjudicate all applicable fields, including uncertain judgments. Report class-specific agreement, audit sampling weights and human estimates separately. This keeps a selectively reviewed item from silently changing the primary endpoint. The adaptive violation rate uses the nine-vote panel described above.

\textbf{Can the frozen model use the rubric directly?} On the full internal fixed test, run the original Llama checkpoint using the pilot's exact target-plus-rubric template, $C(\varnothing,x)$, with the source field empty. It needs only the visible request at inference. This extends XSTest's comparison of the same checkpoint with and without a safety prompt \cite{rottger2023xstest}. On the 250 safe prompts in its \href{https://github.com/paul-rottger/xstest/tree/d7bb5bd738c1fcbc36edd83d5e7d1b71a3e2d84d/model_completions}{released human labels}, Mistral gave a full or partial refusal on 4 prompts without the guardrail and 47 with it. In the paired records, 43 changed to a refusal and none changed in the opposite direction. A recent fixed-item study also found that safety prompts could shift the tendency to refuse \cite{gupta2026systemprompts}.

Join the prompted-base answers to the saved ordinary-base and arm D answers by common test-item ID. Match decoding, output allowance and item weights. Keep overlength items in the denominator as unresolved. Score the visible request and saved answer under the same rubric and criterion, with rendering and checkpoint hidden. Compare prohibited assistance, benign over-refusal, useful benign completion and safe-and-helpful completion on mixed requests. Use the registered family groups and share the two frozen-base observations across D's three seed comparisons. The within-base comparison measures the full prompt effect, including instructions, layout and length. D versus the prompted base compares two deployable systems on these fixed requests. The trained empty-source arm separately tests adding source text to reference scoring on fixed training pairs. Report this secondary comparison's additional generations, grading calls, input and output tokens, and measured serving time.

For the secondary external over-refusal measure, use all 250 prompts labeled safe in the \href{https://github.com/paul-rottger/xstest/blob/d7bb5bd738c1fcbc36edd83d5e7d1b71a3e2d84d/xstest_prompts.csv}{pinned XSTest file} \cite{rottger2023xstest}. The file also contains 200 unsafe contrasts. Exclude these from the over-refusal denominator. Apply the \href{https://github.com/paul-rottger/xstest/blob/d7bb5bd738c1fcbc36edd83d5e7d1b71a3e2d84d/evaluation/classify_completions_gpt.py}{XSTest authors’ classification prompt} with the registered evaluator to each nonempty saved answer. A direct but inaccurate attempt counts as compliance. Offering only an alternative counts as full refusal, and refusing before answering counts as partial refusal. Score usefulness separately. Report partial and full refusals separately and sum them for the any-refusal rate. Keep all 250 IDs in the denominator. If $R$ answers show a refusal and $U$ verdicts remain unresolved after review or retry, report the interval from $R/250$ to $(R+U)/250$ for the observed suite. Flag empty visible answers within $U$ and give them no helpfulness credit. After freezing the XSTest evaluator, check it on the authors' \href{https://github.com/paul-rottger/xstest/tree/d7bb5bd738c1fcbc36edd83d5e7d1b71a3e2d84d/model_completions}{human-labeled completions} for all 250 safe prompts in each of five historical model and prompting configurations. Join by configuration and canonical prompt ID, removing the archived \texttt{v2-} prefix and checking prompt, type and answer text against the \href{https://github.com/paul-rottger/xstest/tree/d7bb5bd738c1fcbc36edd83d5e7d1b71a3e2d84d/evaluation/automated_evaluation_labels}{archived judgments}. Report the three-by-three table of predicted versus adjudicated labels, partial-refusal recall and unresolved predictions separately. Of 123 human-labeled partial refusals among the 1,250 historical answers, the archived GPT evaluator called 17 partial refusals, 82 full compliance and 24 full refusals. These are historical grader results. Keep these prompts and answers out of judge tuning, and report the frozen CRPO grader's own comparison separately from the new models' over-refusal rates. For the primary matched-benign over-refusal rate, use the canonical benign control IDs, including controls linked to several framed variants. A target-only assessor records a full or partial over-refusal when the answer withholds all or part of the requested permissible help on safety grounds without justification under the rubric. A warning followed by all requested help, an inaccurate attempt with no refusal, and a clarification warranted by the request each score zero for over-refusal. Score helpfulness separately. A completed empty reply earns zero helpfulness and leaves the refusal label unresolved. Keep every registered control in the denominator and include unresolved labels in the bounds. For joint accuracy, use the separately registered framed-target/control pairs and require correct handling of both members. Group connected families whenever either scored member of a fixed pair is shared. Adapt the static six-template evaluation to six prespecified, nonadaptive framing styles on a planned 9,160 held-out source goals, matching the scale of the static SEP evaluation. Freeze target texts and validity outcomes before model evaluation and prespecify the 1,024-goal adaptive subset by source family. Report per-style validity, per-style fixed-test violations, and any-of-six success per source including unsuccessful attempts. Prelabel two difficulty dimensions before viewing results: closeness of the prohibited target to its permitted counterpart, and single versus composed framing transformations. Calibrate closeness on the actions requested under the policy. For example, concealment and lawful declaration can involve identical goods. Report ASR and matched benign helpfulness with the family count in each stratum. This adapts SecPO's analysis of semantically close and multiple injections to a single user request. Evaluate the published same-family SecPO and Meta-SecAlign checkpoints on the same fixed framed targets alongside their base. Within each family, the published SecPO and Meta-SecAlign checkpoints share their \texttt{tokenizer.json}, \texttt{tokenizer\_config.json}, chat template and generation configuration (\href{https://huggingface.co/Sizhe-Chen/Llama-3.1-8B-Instruct-MetaSecAlign/tree/dbc4ed4aa435b651c85d47c154a8908fe0f7129f}{Llama 3.1} and \href{https://huggingface.co/Sizhe-Chen/Qwen3.6-27B-Instruct-MetaSecAlign/tree/4a3a6177894ba0d4d73c8d5532aa94ab9adc3076}{Qwen3.6}). Render each standalone request once per defended pair and check that each serving path receives the same input token IDs and decoding settings. Run the undefended base with its registered deployment template. These transfer baselines measure the effect of bringing an existing injection defense to this task. The four matched CRPO/DPO arms control for training data.

\textbf{Does the source improve the reference answer on held-out requests?} Han et al.\ evaluate their frozen teacher with privileged information on held-out inputs \cite[App.~G]{han2026opct}. Apply that comparison to all registered links between a source and a framed target or matched benign control in the fixed test, independently of whether either answer succeeds. Run the frozen Llama base with $C(x^S,x^F)$ and reuse its saved $C(\varnothing,x^F)$ answer. Match the rubric, target, instructions, decoding and output allowance. Use the same fixed-test judge, rubric, visible target and registered criterion to grade each saved answer. Withhold the source, answer origin and model identity. For each applicable outcome, tabulate cases where adding the source corrects an error, introduces one, or leaves the outcome unchanged. For example, a source about smuggling may help the reference withhold concealment advice but cause it to refuse a lawful declaration question.

Join both answers to each registered source--target link. Generate once per distinct rendered input and registered draw. For this diagnostic, divide each target's weight equally among its distinct linked sources and report the number of targets and links by overlap class. Reuse the fixed-test criteria for unsafe assistance, benign over-refusal, useful benign completion and safe-and-helpful handling of mixed requests. Report paired changes and unresolved bounds using the connected family groups. Report this privileged-input result separately from deployment scores. Its observed benefit or harm measures the behavior of the frozen reference on new requests. The trained source-omission control tests the separate question of whether supplying the source during reference scoring improves the deployed student.

\textbf{When a request has both prohibited and permitted parts.} Before inspecting model answers, review every valid full- and partial-overlap fixed target. Record the prohibited assistance the target requests and whether it also asks for a scorable permitted action. A request can retain the entire customs-evasion request and add a question about which declaration form to file. Conversely, dropping one of two prohibited actions can leave a request with no permitted part to answer. Optional suggestions for legal alternatives do not make an item eligible. Record inapplicable and unresolved eligibility separately and freeze the same eligible targets and weights for all checkpoints.

For each eligible target, define from its wording what would fulfill the requested permitted part, including all required elements. Give the same saved answer two judgments: whether it supplies any prohibited assistance under the target policy, and whether it fulfills the permitted part. Report the four possible combinations and bounds for unresolved judgments, overall and by source-overlap class. The secondary safe-and-helpful rate credits an answer only when it withholds all applicable prohibited assistance and fulfills the permitted part. A shared-capability score on one source--target link measures only that capability. For the human audit, reuse the \href{https://arxiv.org/html/2508.09224v1#A1.SS2}{Safe-Completions grading prompt} adapted to the registered permitted part. Distinguish no help, generic safety information, concrete partial help and full useful completion \cite{yuan2025safecompletions}. This applies SecPO's separate task-success and attack-success measurements to one mixed request \cite[Sec.~4.7 and App.~D]{chen2026secpo}.

\textbf{Published fixed framing benchmark.} As a secondary comparison, use AI2's \href{https://github.com/allenai/safety-eval/tree/060cc903d64703214c549b5c3a30ea8ceef2e588/evaluation/tasks/generation/wildjailbreak}{pinned WildJailbreak evaluation files and task configurations} \cite{jiang2024wildteaming}. The files contain 2,000 harmful framings of 200 distinct direct requests and 250 benign framings of 207 distinct direct requests. Report each set separately. Use the supplied HarmBench configuration on the harmful set and the WildGuard benign configuration on the benign set. Both give the model the framed request. The harmful evaluator sees the direct request and the model's answer; the benign evaluator sees the framed request and answer. The \href{https://github.com/allenai/safety-eval/blob/060cc903d64703214c549b5c3a30ea8ceef2e588/evaluation/tasks/generation/base.py#L29-L139}{native scores} report the complement of classified harmful responses and the rate of classified compliance, respectively, with parsing errors counting against each score. Report error counts and denominators. Compliance alone gives no evidence of answer quality, and the harmful grader uses information unavailable to the CRPO fixed-test judge. Preserve the released benchmark's item weights and group variants by their direct request for any uncertainty estimate. These results describe transfer to AI2's benchmark under the named graders. The original paper used a different benign refusal grader.

\textbf{Does framing change the answer?} On the fixed test, pair each full-overlap framed request with its own direct source. Determine eligibility from the texts before scoring the models. Use the same eligible pairs for every checkpoint. Generate both answers with the same deployment decoding settings. Score them separately under the corresponding visible request, rubric and registered shared-capability criterion. Count whether each answer supplied the prohibited capability shared by the two requests. Tabulate the four direct-to-framed outcomes. Let $n_{01}$ count pairs with a violation only on the framed request and $n_{10}$ count pairs with a violation only on the direct request. Report both counts and the net change in violation rate, $(n_{01}-n_{10})/N$, over all $N$ eligible pairs, with bounds for unresolved outcomes. Reuse cached direct and framed answers across all of their registered links. If one answer serves two distinct shared-capability criteria, score it separately against each applicable criterion. Resample connected source families together. Under stochastic decoding the two transitions describe observed draws. Even equally safe contexts can produce discordant answers.

\textbf{Do equally applicable framings receive consistent answers?} A source with one violation in six framings and one with six violations both count as a success in the any-of-six measure. Use OpenSafeIntent's \href{https://github.com/Uppaal/OpenSafeIntent/blob/416ddc0a6655acc48bff9d919e266e20ec81520d/analysis/paraphrase_robustness.py#L66-L115}{released safety-flip classification} on the fixed answers already collected \cite{uppaal2026opensafeintent}. Before viewing answers, group distinct full-overlap target IDs by direct source and shared prohibited-capability criterion. A group needs at least two distinct targets that can be scored against that same criterion. Freeze group membership across checkpoints. For fully scored groups, map withholding the capability to 1 and supplying it to 0. Report the percentages of groups with all answers withholding, a mixture, or all answers supplying the capability. Keep groups with unresolved answers in the denominator. For each percentage report the fraction forced into its category and the fraction that could fall there after completing the missing judgments. These marginal bounds can be conservative when groups share an unresolved judgment. Give group counts by number of variants so that the opportunity to observe a change is visible. With one draw per variant these are observed response patterns; decoding variability can also produce a change.

\textbf{Adaptive test.} Fix 1,024 previously unseen source goals and evaluate each defended model and its base on the same goals. Use two held-out SecPO controllers, TAP and Genetic/MAP-Elites \cite{secpocode}. Route both controllers through the adaptive evaluation wrapper so they share the checkpoint, per-candidate judgment ledger, and authoritative early-stop predicate. Extend the wrapper's witness-required input validation and candidate-judgment call, and replace the substring gate in the SEP vote runner with the confirmed framing-validity decision. A blinded semantic panel then judges every eligible visible answer. Keep the dense critic score solely for parent selection. Preserve model-family stop IDs and response normalization while replacing the wrapper's injection renderer with a single ordinary user message. Give each up to 800 actual defender queries per goal per checkpoint, including seed evaluations. Count every candidate evaluated in a batch and preserve within-batch order for the first verified hit. The released Genetic controller finishes the current generation and TAP finishes expanding the current depth before stopping. Record the first confirmed query and all actual queries spent, including siblings evaluated later in the same depth. Invalid or duplicate rewrites filtered before inference consume proposal effort only. Allow up to four raw proposals per requested candidate slot as a common, provisional retry ceiling. Freeze it on development data. Log underfilled rounds, early exhaustion, and actual queries per goal. For Qwen, the released adaptive preset enables thinking and allows 32,768 output tokens. The main paper describes non-thinking evaluation \cite[Sec.~4.1]{chen2026secpo}. Set \texttt{enable\_thinking=false}, greedy decoding, and an explicit 1,536-token adaptive victim cap for the primary matched run. Pin family stop tokens and pass only the visible final answer to scorers. Record these overrides separately from the AgentDyn serving profile. For Genetic use eight candidates per generation up to 100 generations. For TAP use eight starts, maximum depth ten, branching factor ten and beam width ten, subject to the same cap. The released adaptive configuration defaults to TAP branching factor eight. Record the override to ten for the paper's breadth and width. Preserve Genetic's five islands, capacity 100 per island, diversity archive, and migration every ten generations. Both searches apply the framing validity checks. Report separate curves at 1, 10, 100, and 800 defender queries and a secondary curve indexed by raw proposal count, plus the union of successes explicitly labeled with its larger total budget. For the learned attacker, reuse the \href{https://github.com/pppyb/SecOPD/blob/571502a2a315c4b8820dd878d4569e2a2222cb88/evaluation/pismith/setup.sh}{SecOPD project's pinned PISmith integration} for CRPO. Its setup selects PISmith revision \texttt{16865a7} and applies the Qwen overrides. Train one Qwen3-4B-Instruct-2507 attacker per defender on 100 disjoint development goals. Keep adaptive GRPO, 10 epochs, 15 rollouts per prompt at temperature 1, clipping 0.2, zero KL coefficient and learning rate $10^{-5}$. Use the pinned attacker's 60-completion generation batch (four goals with 15 rollouts each) and its greedy, thinking-off Qwen training victim capped at 1,024 tokens. The cap applies to the victim answer during attacker training. Give a binary reward of one for a valid framing with verified prohibited assistance and zero for a settled unsuccessful attempt. The \href{https://github.com/pppyb/SecOPD/blob/571502a2a315c4b8820dd878d4569e2a2222cb88/evaluation/pismith/overrides/benchmarks/piarena/reward.py}{released reward callback} collapses some service failures into empty answers or negative judgments. Retry unresolved calls on the original sampled group before computing rewards and log any group excluded after retry exhaustion. This matters because adaptive GRPO uses the batch success rate to set its exploration coefficient. At test time reuse the \href{https://github.com/pppyb/SecOPD/blob/571502a2a315c4b8820dd878d4569e2a2222cb88/evaluation/pismith/evaluate.sh}{ten-candidate evaluator}: one batch of ten draws per held-out goal, temperature 0.9 and a 512-token cap per proposal. Preserve candidate order and per-goal pass-at-ten aggregation. The released evaluator fills missing draws with static or repeated text and uses a witness check. For CRPO, preserve genuine draws, including sampled duplicates; retry missing service responses within a fixed retry budget and mark persistent missing slots unresolved. Judge valid framings through the same blinded panel used for TAP and Genetic. Load the preregistered held-out goal manifest directly. The \href{https://github.com/pppyb/SecOPD/blob/571502a2a315c4b8820dd878d4569e2a2222cb88/evaluation/sep/test_sep_adaptive.py}{external PISmith driver} drops rows with an empty SEP data field and accepts a cached batch when its numeric indices cover the requested positions. For CRPO, use permanent goal IDs and ten numbered draw slots per goal. Reuse SecPO's atomic JSON writer and adapt its checkpoint fingerprint to include the ordered manifest, frozen attacker and defender revisions, rendering, decoding, and judge configuration. Require exact slot coverage in completed results; a reordered cohort or changed checkpoint requires a fresh evaluation. The \href{https://huggingface.co/AlbertYin/piarena_attacker_qwen3_4b_secalign/tree/a34665c4b98de79245d4e7bd046a43719408c124}{public PIArena attacker} supplies a native serving fixture. Report learned-attacker ASR@10 and training cost from the attacker trained for each defender. Define cumulative confirmed success at budget $b$ as
\[
\mathrm{ASR}@b =
\frac{\#\{\text{held-out source goals with a verified failure within }b\text{ defender queries}\}}
     {\#\{\text{held-out source goals attempted}\}}.
\]
SecPO's SEP panel judges whether the answer carried out the particular injected instruction \cite[Sec.~4.1 and App.~B.4]{chen2026secpo}. CRPO uses the same interpretation for its primary adaptive rate. A source goal counts once if the panel confirms that an answer to a valid framing delivered the prohibited capability recorded for that source--target pair within budget. On a partial-overlap framing, the verdict concerns the retained part. An answer that withholds a requested private record but discloses a different one is a confidentiality violation and should be recorded separately. It contributes zero to success on the requested record. The fixed direct-to-framed comparison reports its own rate over eligible full-overlap pairs.

Before aggregation, reconcile one terminal search record per prespecified goal, attack, checkpoint and seed against the frozen manifest. Adapt the completeness checks used by SecPO's AgentDojo and DTaP evaluators \cite[App.~B.5]{chen2026secpo}. A search that exhausts its allowed proposals with no valid framing is a completed attempt with zero confirmed successes. A missing run, a stale checkpoint result, or a crash before the stopping condition leaves an incomplete run. Resume it before computing the primary comparison. A completed search with exhausted judge retries retains its recorded uncertainty. Show coverage and provisional bounds if execution remains incomplete. Every allocated goal enters the final denominator. Report validity yield, unresolved judgments, defender queries and total generation/judging calls separately. Among raw proposals, count explicit proposer refusals separately from malformed, duplicate, and policy-invalid rewrites. This shows how often the attacker declined to generate a candidate before the defender had an opportunity to answer \cite[Sec.~7]{chen2026secpo}. Plot success against budget and record the first verified success per goal. Use fresh answers on a subset to assess repeatability of discovered failures. Report unconditional cost per attempted goal. Among successful goals, give mean and median defender queries and cost, with the number of successes, token counts, and pricing date \cite[App.~C.7]{chen2026secpo}. The same held-out goals allow paired comparisons even when each defender elicits different framings.

\textbf{Final judging and inference.} Within each TAP or Genetic comparison, use the same proposer checkpoint or service revision, decoding settings, prompt template, validity judges, search critic, calibration examples, and final panel across defenders. Put the proposer and critic identities in the SEP runner's \texttt{AdaptiveExample.metadata} before search. Its existing checkpoint check covers this field. Reuse the DTaP \href{https://github.com/Sizhe-Chen/SecPO/blob/da0372a0e598badf4e51c98c175a962147cac2ac/src/secpo/evaluation/dtap_qwen36.py}{\texttt{backend\_fingerprint}} helper on the callback functions, recording the configured model routes, prompt hashes and decoding settings through its runtime-identity hook \cite{secpocode}. A changed backend then starts a separate search history. The learned PISmith attacker remains defender-specific by design. Give adaptive searches common initial candidates for each goal and let each defender's answers determine subsequent proposals. Keep model-family-specific preference judges restricted to training. Freeze the final rubric and decisions about material safety and helpfulness regressions before looking at held-out outputs. Final assessors see the target, policy, and anonymized answer, with search history, scores, and model identity withheld. A separate source-aware check verifies only the source--target relation. The common frozen judge supplies the primary internal fixed-test labels. On adaptive tests the panel supplies the registered verdict. Blind human review covers predicted successes, unsuccessful attempts, benign refusals, and disagreements across arms. Report adjudicated human judgments and their agreement with the registered scorers by class. Repeated calls to one checkpoint share its biases. If unresolved cases could reverse a comparison, report the corresponding range. Estimate uncertainty by resampling whole source families together across arms. For marginal binary rates on source goals drawn independently from one common population, report two-sided 95\% Clopper--Pearson intervals. Use \href{https://github.com/scipy/scipy/blob/v1.18.1/scipy/stats/_binomtest.py}{SciPy 1.18.1's \texttt{binomtest}} with \texttt{proportion\_ci(method='exact')}. Two published base-model results check the implementation: 47 of 198 correct on Llama GPQA Diamond gives 18.0--30.3\%, and 87 of 97 completed clean Qwen AgentDojo tasks gives 81.9--94.9\% \cite[Tables~13 and~15]{chen2026secpo}. For prespecified domain mixtures, use the stratified procedure below. Retain paired family-level inference for model differences and separately report training-seed variation.

\subsection{General utility and the decision rule}
Evaluate each base and trained checkpoint on clean tasks in its normal deployment configuration. Freeze prompts, model revisions, generation limits, and scoring procedures before final evaluation. Report the absolute score and paired change from that model family's own base for each benchmark. In SecPO's evaluation harness, the following settings measure different capabilities \cite[App.~B.3]{chen2026secpo}.

\begin{itemize}
\item \textbf{Knowledge and reasoning.} Run the full MMLU-Pro (five-shot), GPQA Diamond (zero-shot), GSM8K (eight-shot), and Minerva Math (four-shot) suites. Record the task definitions and harness revision actually used. Use the \href{https://github.com/Sizhe-Chen/SecPO/blob/da0372a0e598badf4e51c98c175a962147cac2ac/src/secpo/evaluation/lm_eval_utility.py}{released utility launcher} for the four tasks. Load a PEFT adapter on its exact base via \texttt{--adapter-path}. A merged or otherwise complete checkpoint, including each Qwen3.8 export, is itself the \texttt{pretrained} model. Supply its own model ID and revision and omit the adapter argument. For a validated local export, set \texttt{pretrained} to the local directory in the generated harness command and omit the Hub revision for those local weights. Save the final command and export-manifest hash, and verify which weights and PEFT state the harness loaded. Keep the same tasks, tokenizer, template and decoding in matched comparisons. For the main Qwen3.6 and secondary thinking-off Qwen3.8 utility runs, add \texttt{enable\_thinking=false} to the generated \texttt{--model\_args} value for every checkpoint in that comparison. The released launcher sets chat formatting but omits this argument. The pinned \href{https://github.com/EleutherAI/lm-evaluation-harness/blob/d6de81643928d653435c431bae19945d41d32520/lm_eval/models/huggingface.py}{harness} leaves an omitted value unset, and the \href{https://huggingface.co/Qwen/Qwen3.6-27B/blob/6a9e13bd6fc8f0983b9b99948120bc37f49c13e9/chat_template.jinja}{Qwen template} then opens a reasoning block. Its Boolean false branch closes the block before generation. Pass false as an unquoted Boolean in the model arguments. The string containing the word false opens the block. Check the effective Boolean and rendered assistant prefix on one item from each of the four tasks. Record the command adjustment alongside the native dry run. The release supplies \href{https://github.com/Sizhe-Chen/SecPO/blob/da0372a0e598badf4e51c98c175a962147cac2ac/configs/evaluate_lm_eval_llama31_8b.json}{launcher configurations}; its \href{https://github.com/Sizhe-Chen/SecPO/blob/da0372a0e598badf4e51c98c175a962147cac2ac/docs/EVALUATION.md#utility-evaluation}{evaluation notes} list the original custom Llama task files as external \cite{secpocode}. For a direct paper comparison, recover those definitions and pin their hashes with the harness revision. In the existing utility smoke test, inspect a rendered prompt and a scored fixture per task. Verify the \href{https://arxiv.org/html/2610.41998v1#A2.SS3}{paper's chain-of-thought instruction} and total demonstration count, including examples embedded in the task prompt \cite[App.~B.3]{chen2026secpo}. Check the extracted choice, number, or mathematical answer against the expected fixture result. The public SecPO repository and Llama model and dataset releases inspected do not identify the original task files. For the reconstruction, use the \href{https://github.com/EleutherAI/lm-evaluation-harness/tree/d6de81643928d653435c431bae19945d41d32520/lm_eval/tasks}{EleutherAI task definitions at \texttt{d6de816}} with SecPO's four task names, extraction metrics and full splits. For MMLU-Pro, read \texttt{exact\_match,custom-extract} from the \href{https://github.com/EleutherAI/lm-evaluation-harness/blob/d6de81643928d653435c431bae19945d41d32520/lm_eval/tasks/mmlu_pro/_mmlu_pro.yaml}{group} and \href{https://github.com/EleutherAI/lm-evaluation-harness/blob/d6de81643928d653435c431bae19945d41d32520/lm_eval/tasks/mmlu_pro/_default_template_yaml}{five-shot template}. For \href{https://github.com/EleutherAI/lm-evaluation-harness/blob/d6de81643928d653435c431bae19945d41d32520/lm_eval/tasks/gpqa/cot_zeroshot/gpqa_diamond_cot_zeroshot.yaml}{GPQA Diamond} with its \href{https://github.com/EleutherAI/lm-evaluation-harness/blob/d6de81643928d653435c431bae19945d41d32520/lm_eval/tasks/gpqa/cot_zeroshot/_gpqa_cot_zeroshot_yaml}{zero-shot template}, and for \href{https://github.com/EleutherAI/lm-evaluation-harness/blob/d6de81643928d653435c431bae19945d41d32520/lm_eval/tasks/gsm8k/gsm8k-cot.yaml}{eight-shot GSM8K}, use \texttt{exact\_match,flexible-extract}. For the \href{https://github.com/EleutherAI/lm-evaluation-harness/blob/d6de81643928d653435c431bae19945d41d32520/lm_eval/tasks/benchmarks/minerva_math.yaml}{seven-subject Minerva MATH group}, read \texttt{math\_verify,none} from each \href{https://github.com/EleutherAI/lm-evaluation-harness/blob/d6de81643928d653435c431bae19945d41d32520/lm_eval/tasks/minerva_math/minerva_math_algebra.yaml}{subject task} and weight subjects by their number of test items. Apply them identically to the base, DPO and CRPO runs. The Meta-SecAlign \href{https://github.com/facebookresearch/Meta_SecAlign/blob/2031502bdf5a046b437ec835c6b05ff468aa0699/lm_eval_config/mmlu_pro/mmlu_pro_5shot_cot_instruct.yaml}{MMLU-Pro} and \href{https://github.com/facebookresearch/Meta_SecAlign/blob/2031502bdf5a046b437ec835c6b05ff468aa0699/lm_eval_config/gpqa_cot/gpqa_0shot_cot.yaml}{GPQA} definitions read Meta's preformatted prompts and use a strict-match filter. The \href{https://github.com/pppyb/SecOPD/blob/571502a2a315c4b8820dd878d4569e2a2222cb88/evaluation/lm_eval/run_paper_suite.sh}{SecOPD launcher} runs full GPQA and GSM8K but samples 100 items per category for MMLU-Pro and Minerva. CRPO uses the full splits. Label reconstructed absolute scores separately from Table~1. Keep the seed and chat template fixed within each comparison. Give each checkpoint its own output directory. Read each task's configured metric key and join its logged samples by benchmark item for paired changes. Count answer-extraction failures using the benchmark scorer. For GPQA Diamond the \href{https://github.com/EleutherAI/lm-evaluation-harness/blob/d6de81643928d653435c431bae19945d41d32520/lm_eval/tasks/gpqa/cot_zeroshot/utils.py}{pinned harness} shuffles the answer options and derives the correct letter from that order. Freeze the processed options and key once, and match the rendered question, option order and key per item across checkpoints within a family. Reuse the item-level check in the \href{https://github.com/0315jaewon/SecPO-test/blob/42280080f0f9224717f40a53928a761e554e1db0/scripts/lambda_gpqa.py#L239-L282}{coauthor GPQA scoring replay}. From each CRPO run’s saved responses, require each of the 198 document IDs exactly once per registered filter, then reapply it and verify the extracted answer and exact-match score against the saved values. The original launcher expects a particular Qwen3.8 checkpoint and historical run files. Chang's separate \href{https://github.com/0315jaewon/SecPO-test/blob/d009be8d9779f05cd475b3e7a490e829fff12d64/docs/LAMBDA_GPQA.md}{GPQA release} has the same generation and scoring script as the archive and pins separate environments for the two steps. For the historical Qwen3.8 base replay, the later \href{https://github.com/0315jaewon/SecPO-test/blob/d009be8d9779f05cd475b3e7a490e829fff12d64/scripts/export_lambda_gpqa_bundle.py}{input exporter} can be used unchanged when the original private reference files are available. It validates the frozen requests, reference records and five tokenizer files and writes per-file checksums. Its synthetic tests exercise those checks without the benchmark data. For new CRPO evaluations, freeze the corresponding files for the registered checkpoint and retain the item-level scoring checks above. For MMLU-Pro and Minerva Math, reuse the coauthor runner’s \href{https://github.com/0315jaewon/SecPO-test/blob/42280080f0f9224717f40a53928a761e554e1db0/scripts/evaluate_lambda_full_lm.py}{full-split checks} against the registered CRPO task definitions. Before scoring, verify the split size, dataset fingerprint and selected document contents. Before aggregation, require every declared scoring filter and match its ordered item IDs and document contents within each subtask.
\item \textbf{Open-ended instruction following.} Run all \href{https://github.com/Sizhe-Chen/SecPO/blob/da0372a0e598badf4e51c98c175a962147cac2ac/configs/evaluate_alpacaeval2_utility.json}{805 AlpacaEval2 instructions} and all \href{https://github.com/Sizhe-Chen/SecPO/blob/da0372a0e598badf4e51c98c175a962147cac2ac/configs/evaluate_sep_utility.json}{9,160 clean SEP items} with their pinned reference responses. Use SecPO's raw win rate against the archived GPT-4o-2024-05-13 answers for AlpacaEval2 and Meta-Llama-3-8B-Instruct answers for SEP, with Claude-Opus-4.8 as the paper's evaluator. The \href{https://github.com/tatsu-lab/alpaca_eval/blob/cd543a149df89434d8a54582c0151c0b945c3d20/results/gpt-4o-2024-05-13/model_outputs.json}{public GPT-4o answer file} contains 805 unique instructions labeled \texttt{gpt-4o-2024-05-13} and passes SecPO's unmodified reference preflight. The \href{https://github.com/egozverev/Should-It-Be-Executed-Or-Processed/blob/7606c0696f20f5aa433169fd2221f76852d1d4f5/SEP_dataset/SEP_dataset.json}{original 9,160 SEP records} supply the underlying questions and witnesses. For SEP, use the fixed \href{https://github.com/0315jaewon/SecAlign_cont/blob/9a1ae20a6421a834d46444fbe35d699e2bcfaacc/data/SEP_dataset_test_Meta-Llama-3-8B-Instruct.json}{Meta-Llama-3-8B-Instruct answers} and \href{https://github.com/0315jaewon/SecAlign_cont/blob/9a1ae20a6421a834d46444fbe35d699e2bcfaacc/data/SEP_dataset_test.json}{ordered test manifest} published by SecPO coauthor Jaewon Chang. These files pass the \href{https://github.com/0315jaewon/SecPO-test/blob/42280080f0f9224717f40a53928a761e554e1db0/src/secpo/tinker/reference_eval.py#L69-L106}{authors’ alternate-reference validator} on all 9,160 rows. It verifies the processed manifest digest and substitutes the alternate answer digest in memory. It checks every ordered instruction, input, injection and witness, the nonempty answer, and the generator \texttt{meta-llama/Meta-Llama-3-8B-Instruct}. Supply its returned configuration and the candidate file to the \href{https://github.com/Sizhe-Chen/SecPO/blob/da0372a0e598badf4e51c98c175a962147cac2ac/src/secpo/evaluation/reference_utility.py}{SecPO preflight and command builder}. The exact manuscript comparator remains unverified. Use the same alternate across all arms, label its scores \emph{SEP utility, public alternate reference}, and record the following SHA-256 digests. The last row gives the expected manuscript digest.
\begin{center}
\scriptsize\begin{tabular}{@{}ll@{}}
\textbf{File} & \textbf{SHA-256} \\
AlpacaEval2 answers & \texttt{653c131d58e8d5802203c85e80ca3fb4b378f22e2a0280ad5cd910ff0bee9535} \\
SEP alternate answers & \texttt{63e965d2a25d2d561a7b6079048977314cfaacf5d286cf86ae168538c771dd34} \\
SEP processed manifest & \texttt{e659683ae5b10905064d28bca5fd786d1a4ec96ce016a6b93a7282fd977a95a3} \\
SEP manuscript answers & \texttt{6defcfdb3ac9da4a03eff21403a0ac58ce52c0c1dc45c23fc48b17a2bba0808c} \\
\end{tabular}
\end{center}
 Use Chang's \href{https://github.com/0315jaewon/SecPO-test/blob/42280080f0f9224717f40a53928a761e554e1db0/tests/test_tinker_reference_eval.py#L45-L65}{clean-input test} to check that only \texttt{instruction\_only} and optional \texttt{input} reach the candidate through its deployment template. For each within-family comparison, reuse the \href{https://github.com/0315jaewon/SecPO-test/blob/42280080f0f9224717f40a53928a761e554e1db0/scripts/evaluate_tinker_reference.py#L73-L158}{evaluation runner's protocol check and row expansion}. Give each base and trained checkpoint an explicit output directory. Compare the ordered rendered-prompt digest, references, tokenizer, reasoning mode, decoding, judge, metric and code revisions. A resumed run must also match the checkpoint identity. Only that identity may differ between paired runs. Keep a separate cache per checkpoint and protocol. Generate once per distinct prompt under deterministic decoding and restore the candidate to every original row. Score all 9,160 rows against their own reference answers and cluster paired comparisons by the 9,135 distinct clean prompts. Three of the 23 repeated prompts have different reference answers across rows. Hold the judge configuration fixed across arms and record its checksum and any substitution. Before aggregation, normalize a documented legacy tie of 0 to 1.5. The test fork's \href{https://github.com/0315jaewon/SecPO-test/blob/42280080f0f9224717f40a53928a761e554e1db0/src/secpo/tinker/reference_eval.py#L282}{\texttt{validate\_annotations}} requires a finite preference in $[1,2]$ and reconciles every judgment with the ordered item, both answers and both generator identities. Its \href{https://github.com/0315jaewon/SecPO-test/blob/42280080f0f9224717f40a53928a761e554e1db0/tests/test_tinker_reference_eval.py#L121}{answer-order fixture} verifies candidate-relative scoring under reversed presentation and excludes SEP injection metadata from the clean judge prompt. For both suites, read \texttt{win\_rate} from AlpacaEval's \href{https://github.com/tatsu-lab/alpaca_eval/blob/cd543a149df89434d8a54582c0151c0b945c3d20/src/alpaca_eval/metrics/winrate.py}{\texttt{get\_winrate}} using the archived reference as answer 1 and the candidate as answer 2. The \href{https://github.com/tatsu-lab/alpaca_eval/blob/cd543a149df89434d8a54582c0151c0b945c3d20/src/alpaca_eval/metrics/helpers.py}{raw scoring rule} averages $p-1$ over valid preferences $p\in[1,2]$ and multiplies by 100. A tie contributes 0.5 and a preference of 1.9 contributes 0.9. A fixture containing just those two ratings must score 70 with two valid judgments. Publish \texttt{n\_total} with the raw win rate, label incomplete scores with their coverage, and retry unresolved rows on the saved answers. A full-suite comparison requires 805 valid judgments for AlpacaEval2 and 9,160 for SEP per checkpoint. On an interrupted utility run, use Chang's \href{https://github.com/0315jaewon/SecPO-test/blob/42280080f0f9224717f40a53928a761e554e1db0/scripts/judge_lambda_cached_reference.py#L37-L82}{saved-answer validator and cache-only sampler} when the serving records use its Qwen schema. Check the ordered rendered prompts, checkpoint and decoding identity, candidate text, generation records and token counts against the saved completion cache. A missing or mismatched entry blocks that row from resumed judging. For another serving format, bind its actual stop rules and request identity to the cache and exercise the \href{https://github.com/0315jaewon/SecPO-test/blob/42280080f0f9224717f40a53928a761e554e1db0/tests/test_lambda_judging_resume.py#L22-L76}{duplicate-prompt and changed-answer fixtures}. Keep candidate generation fixed while retrying judgments. Use the \href{https://github.com/0315jaewon/SecPO-test/blob/42280080f0f9224717f40a53928a761e554e1db0/tests/test_tinker_evaluation_variants.py}{checkpoint-comparison tests} to check cache isolation and that only the registered weights and associated routing differ. If the judge changes, score every compared arm's saved answers under the same version and report that comparison separately. Chang's \href{https://github.com/0315jaewon/SecPO-test/blob/42280080f0f9224717f40a53928a761e554e1db0/scripts/analyze_lambda_judge_shift.py#L47-L99}{historical judge replay} can check parser and presentation-order effects when the original judge replies and its supported response schema are available. Record answer lengths. For clean SEP, render each instruction and its data in the same ordinary deployment format for every checkpoint of that family. Inspect newly introduced refusals, irrelevant warnings, and incomplete answers.
For AlpacaEval2, also report a secondary length-controlled win rate alongside the raw score and each checkpoint's change from its own base. This checks whether the utility comparison is sensitive to a judge's preference for longer answers. Reuse Chang's \href{https://github.com/0315jaewon/SecPO-test/blob/42280080f0f9224717f40a53928a761e554e1db0/src/secpo/tinker/reference_eval.py#L282-L319}{annotation validator and scorer} on the same 805 judgments in canonical instruction order. The \href{https://github.com/tatsu-lab/alpaca_eval/blob/cd543a149df89434d8a54582c0151c0b945c3d20/src/alpaca_eval/metrics/glm_winrate.py}{AlpacaEval 0.6.6 implementation} joins instruction difficulty by row position. Verify the instructions against the ordered manifest, pass records in positions 0 through 804, and check the \href{https://huggingface.co/datasets/tatsu-lab/alpaca_eval/blob/2edc6fad8be6b14ea7230aabfd08188da6b8b814/df_gamed.csv}{calibration file} against the scorer's pinned digest. The fitted model estimates the candidate's win rate with the character-length difference set to zero \cite{dubois2024length}. Report the adjusted score as unavailable if the judgments are incomplete, the length differences have zero variance, or the fit fails to return a finite score. The adjusted estimate is descriptive for this judge and comparator. SecPO's raw scores determine the primary preservation test and six-suite average.

\item \textbf{Benign tool use.} Run AgentDojo's 97 unattacked tasks on the base and trained checkpoints using identical tool schemas, chat wrappers, tool parser, and environment version within a family. For the paper-compatible score, pool successes across all 97 registered tasks. The \href{https://github.com/Sizhe-Chen/SecPO/blob/da0372a0e598badf4e51c98c175a962147cac2ac/src/secpo/evaluation/agentdojo_qwen36.py#L1174-L1185}{upstream summary} averages four suite rates equally, although its suites contain 16, 21, 20 and 40 tasks. Report this suite average separately. SecPO Table~15's utility rates and exact binomial intervals agree with 87, 85 and 88 successes out of 97 for Qwen3.6 base, Meta-SecAlign and SecPO, respectively. Use task-weighted completion for the figure and report tool-format errors separately. Add all 60 clean AgentDyn tasks for Qwen using the \href{https://github.com/Sizhe-Chen/SecPO/blob/da0372a0e598badf4e51c98c175a962147cac2ac/docs/AGENTDYN.md}{pinned adapter and upstream clean runner}. For Qwen3.8, reuse Chang's \href{https://github.com/0315jaewon/SecPO-test/blob/42280080f0f9224717f40a53928a761e554e1db0/scripts/evaluate_tinker_agents.py}{clean episode runner} and \href{https://github.com/0315jaewon/SecPO-test/blob/42280080f0f9224717f40a53928a761e554e1db0/src/secpo/tinker/agent_eval.py}{schema-aware tool parser and serializer}. Register the new checkpoint, renderer and decoding settings. Preserve task IDs and resume individual episodes under the identical configuration. Keep unresolved infrastructure errors visible in coverage and bounds on the full task set. Record task completion, tool-format errors, and integration failures. Also review the five clean AgentDyn cases in \href{https://arxiv.org/html/2610.41998v1#A4.T21}{SecPO Table~21} \cite[App.~D]{chen2026secpo} using the existing transcripts, including when both checkpoints receive the same completion score. For each, record the fact returned by the tool and whether the subsequent action satisfies the user’s condition in light of that fact. The five cases are \texttt{shopping/user\_task\_16}, \texttt{shopping/user\_task\_14}, \texttt{dailylife/user\_task\_5}, \texttt{dailylife/user\_task\_19} and \texttt{github/user\_task\_3}. In the second shopping case, the rating and price determine whether the requested purchase is appropriate. These cases illustrate behavior within the 60-task suite. For clean tasks whose completion changes between a base and a tuned checkpoint, inspect the anonymized transcripts. Record unwarranted refusals, warranted requests for clarification, tool-call errors and other causes, allowing uncertainty. Report gains and losses against the common task counts alongside the native completion score. The \href{https://github.com/ethz-spylab/agentdojo/blob/d3640b5b03a88eb44dad96852e1de1ef437d836e/src/agentdojo/benchmark.py#L278-L301}{pinned clean runner} assigns zero utility on certain context-limit and server errors. Keep that zero in the native 97-task score and flag its cause. When a terminal error follows a tool action and the partial environment and trace were captured, apply the \href{https://github.com/chawins/agentdojo/blob/07e5ed5d6039571eb7b5bda683143ebd2666e1b7/src/agentdojo/task_suite/task_suite.py#L433-L477}{coauthor fork’s task checker on partial state} to determine whether the requested action had already succeeded. Report this diagnostic separately and classify cases without recoverable evidence as unresolved in the error analysis. Check that calculation on the \href{https://github.com/chawins/agentdojo/tree/1594475cb78cb2c0ce2c8ea8cb75578fce3d9490/runs}{archived clean GPT-4o runs in Chawin Sitawarin's AgentDojo fork}. Join ordinary and \texttt{repeat\_user\_prompt} records by suite and user task, selecting records without an injection task or attack. Across \href{https://github.com/ethz-spylab/agentdojo/tree/d3640b5b03a88eb44dad96852e1de1ef437d836e/runs}{all 97 tasks, 66 succeed under both settings}, 16 gain, one loses and 14 fail under both. Task-weighted completion rises from 67/97 (69.07\%) to 82/97 (84.54\%). On those same runs, the unweighted suite average rises from 70.86\% to 86.18\%. These counts validate the scoring procedure before evaluating CRPO checkpoints. SecPO describes \href{https://arxiv.org/html/2610.41998v1#S7}{excessive refusals in a frontier model}'s clean AgentDyn transcripts \cite[Sec.~7]{chen2026secpo}. Reuse SecPO's \href{https://github.com/Sizhe-Chen/SecPO/blob/da0372a0e598badf4e51c98c175a962147cac2ac/src/secpo/evaluation/agentdojo_qwen36.py}{checkout verifiers} before comparing checkpoints and record their receipts with the utility scores. AgentDojo accepts its pinned \href{https://github.com/Sizhe-Chen/SecPO/blob/da0372a0e598badf4e51c98c175a962147cac2ac/docs/AGENTDOJO.md}{pristine tree} or the exact optional resumption patch. AgentDyn verifies the patched benchmark tree and the Qwen template. On hosts with a symlinked temporary path (as observed under \texttt{/var/folders} on macOS), run the AgentDyn verifier with a canonical temporary directory so Git can apply the patch to an independent copy. On macOS, \texttt{TMPDIR=/private/tmp} reproduced patched-tree SHA-256 \texttt{c0cc7b493d3239c5\allowbreak{}7e7de622f398a8e0\allowbreak{}4f8b6f102b690ca3\allowbreak{}1d99f43695fb353f} and the verifier rejected an additional file edit. Its released serving template enables thinking. Report the released serving mode separately from the matched non-thinking base/CRPO condition, and pin the mode for both sides of each comparison. For Qwen on AgentDojo, also record \href{https://github.com/Sizhe-Chen/SecPO/blob/da0372a0e598badf4e51c98c175a962147cac2ac/configs/evaluate_agentdojo_qwen36_27b.json}{benchmark version v1.2.1}, the \texttt{repeat\_user\_prompt} defense, and the exact task checker. The \href{https://huggingface.co/Sizhe-Chen/Qwen3.6-27B-Instruct-SecPO/blob/4009c2799ed23c4cbb01fd49883c96fe2f2b3dd2/chat_template.jinja}{November 10 Qwen3.6 SecPO standalone template} serializes tool calls with function and parameter tags and renders returned tool messages in the input role. The \href{https://huggingface.co/Sizhe-Chen/Qwen3.6-27B-Instruct-SecPO/blob/4009c2799ed23c4cbb01fd49883c96fe2f2b3dd2/tokenizer_config.json}{template embedded in the same tokenizer configuration} still uses JSON calls and renders those messages in the user role. The November 10 \href{https://huggingface.co/Sizhe-Chen/Qwen3.6-27B-Instruct-MetaSecAlign/blob/d850ad5cda5155e291af5e091b67bb5f6c7f233f/chat_template.jinja}{Qwen3.6 Meta-SecAlign template} has the same split. Both standalone hashes begin \texttt{479393bb39d1}, and both embedded hashes begin \texttt{fc7c6fc019f4}. Before clean tool-use scoring, replay a benign lookup and its result through the actual wrapper and parser. Record the recovered function and argument, rendered result role, and template hash. Use the published benchmark wrapper for its historical control and the same deployment renderer across paired CRPO comparisons. The \href{https://github.com/Sizhe-Chen/SecPO/blob/da0372a0e598badf4e51c98c175a962147cac2ac/src/secpo/evaluation/agentdojo_models.py}{released Qwen AgentDojo wrapper} enables thinking and preserves it across tool calls. Run its native configuration for the upstream control and freeze a common serving mode for each paired CRPO/base comparison. Mark the Llama wrapper as a project integration. Use the \href{https://github.com/Sizhe-Chen/SecPO/blob/da0372a0e598badf4e51c98c175a962147cac2ac/scripts/merge_qwen36_27b_lora.py}{released Qwen LoRA merge path} for external harnesses and verify that the served merged model agrees with adapter-enabled inference on a fixed set of examples before benchmarking. In a separate Qwen3.6 AgentDojo study, Sitawarin's group \href{https://github.com/chawins/agentdojo/blob/07e5ed5d6039571eb7b5bda683143ebd2666e1b7/runs/SUMMARY.md}{traced a base-identical export to unmatched adapter keys}. Reuse its \href{https://github.com/chawins/agentdojo/blob/07e5ed5d6039571eb7b5bda683143ebd2666e1b7/scripts/qwen3_6_eval/merge_lora.py#L248-L283}{per-key loading check and targeted-weight comparison} with the SecPO merger. Reconcile each expected LoRA A and B tensor with its destination after key mapping. A missing or unexpected adapter key invalidates the export even when other layers give a nonzero combined norm. For a targeted weight whose predicted update survives BF16 rounding, compare the saved values with the base plus that update at the deployment precision. If no targeted weight has a representable change, record the export as base-equivalent at that precision.
\item \textbf{Clean document understanding (Qwen).} Test the BF16 base and all four offline Qwen adapters on the same unaltered document images and questions. \href{https://arxiv.org/html/2610.41998v1#S4.SS5}{SecPO Table~5} uses benign DocVQA inputs from Repeat-After-Me, which selects the \href{https://arxiv.org/html/2609.04533v1}{first 100 unique validation images} \cite[Sec.~4.5]{chen2026secpo}\,\cite{chen2026repeat}. Use the \href{https://github.com/Sizhe-Chen/SecPO/blob/da0372a0e598badf4e51c98c175a962147cac2ac/src/secpo/model_merge_qwen36.py}{vision-capable merged checkpoint} and the image processor from the pinned base revision. The \href{https://huggingface.co/Sizhe-Chen/Qwen3.6-27B-Instruct-SecPO/blob/726da6ba0f153d478c207dd5edc593e18e7fd513/merge_manifest.json}{published Qwen merge record} supplies a concrete preflight for the native control. Its index lists 1,184 tensors, including 333 vision tensors, and its \texttt{preprocessor\_config.json} matches the pinned base at SHA-256 {\small\texttt{2722\allowbreak{}5450\allowbreak{}ac9c\allowbreak{}6529\allowbreak{}872e\allowbreak{}e192\allowbreak{}4fcb\allowbreak{}0962\allowbreak{}ff56\allowbreak{}3483\allowbreak{}4f81\allowbreak{}7040\allowbreak{}f444\allowbreak{}1181\allowbreak{}16f4\allowbreak{}e516}}. The authors subsequently \href{https://huggingface.co/Sizhe-Chen/Qwen3.6-27B-Instruct-SecPO/commit/51bd0e76a741defb749be4a973f20d1c1d6ada5c}{corrected the vision tensor names for vLLM}. All 333 moved from \texttt{model.language\_model.visual.} to \texttt{model.visual.}. Both releases have the same counts, merge manifest and image processor. For a new Qwen3.6 export with the nested keys, use Sitawarin's \href{https://github.com/chawins/agentdojo/blob/07e5ed5d6039571eb7b5bda683143ebd2666e1b7/scripts/qwen3_6_eval/merge_lora.py#L75-L136}{vision-weight remapper} on a staging copy and its \href{https://github.com/chawins/agentdojo/blob/07e5ed5d6039571eb7b5bda683143ebd2666e1b7/tests/scripts/test_qwen3_6_eval.py#L178-L279}{shard and index fixtures}. Check index-to-shard agreement and untouched text weights before promoting the export. Applying the prefix change reproduces the corrected index including every shard assignment. The base and SecPO indexes have different total key counts. For each BF16 Qwen3.6 DPO or CRPO export served through vLLM, match the full set of vision keys to the corrected index, check the actual shard headers with the released shard checker, and verify the pinned processor checksum. Then replay the clean color check using the \href{https://huggingface.co/Sizhe-Chen/Qwen3.6-27B-Instruct-SecPO-8Bit/blob/8674c489dda144c51441b5a75cfb7cfd2f7d49f8/vision_smoke_red.png}{published image} and the question in the \href{https://huggingface.co/Sizhe-Chen/Qwen3.6-27B-Instruct-SecPO-8Bit/blob/8674c489dda144c51441b5a75cfb7cfd2f7d49f8/validation_results.json}{validation record}. The authors recorded the visible answer ``Red'' after reloading the bitsandbytes INT8 checkpoint. Preserve the recorded roles for that native control. Pass the same image and question through the ordinary image-and-text template for the BF16 Qwen3.6 base and the four merged DPO and CRPO models evaluated on DocVQA. Log the visible answer as a serving check separately from the benchmark score. Keep preprocessing, chat template and decoding fixed across runs. The authors identify the Table~5 metric as validation average normalized Levenshtein similarity (ANLS) in their \href{https://huggingface.co/Sizhe-Chen/Qwen3.6-27B-Instruct-SecPO/blob/28c89b1e4c28aa47cd048f9c0f7616ad636cf028/README.md}{archived model card}. The \href{https://huggingface.co/Sizhe-Chen/Qwen3.6-27B-Instruct-SecPO/blob/017eb59a72fe06f5683c3bae6c834849a810b3f3/README.md}{earlier model card} labels the same scores of 96.0, 95.8 and 95.2 (base, Meta-SecAlign and SecPO) as thinking-on. The later card describes them as thinking-off. Fix thinking off for every CRPO DocVQA comparison and verify it in the rendered multimodal prompt. The historical mode, exact question IDs, prompts and scorer remain unverified. If those artifacts remain unavailable, freeze the validation revision and row order, take the first question for each of the first 100 unique images, and use the \href{https://github.com/google-research/pix2struct/blob/110d6666e8813936185fa0e9e6c38f35b387012c/pix2struct/metrics.py}{Pix2Struct ANLS implementation} \cite{mathew2021docvqa}. It lowercases the strings, takes the best score across reference answers per question, and averages on a 0--100 scale. Check an exact match, a match to a second reference answer, and a normalized edit distance of exactly 0.5, which earns zero credit. Record item selection, answer extraction and scorer revision as the reconstruction. Report the absolute score and paired change from base by image as a secondary utility endpoint, outside the six-suite text mean.
\end{itemize}

\textbf{Does tuning change predictions on clean answers the base model already gets right?} Use the measurement in RAISED, Eq.~2 and App.~B.1 \cite{dhouib2026raised}, as a secondary diagnostic on Llama. Assemble a held-out set of 125 questions each from ARC-Easy, TriviaQA, GSM8K and SciQ. Register the dataset revisions, prompts, answer checks and selection order. Screen out questions likely to elicit a refusal or guarded reply, then retain items whose greedy base answer passes the registered correctness check. Freeze the item IDs and base answer tokens before inspecting any trained model, and report realized coverage. For each paired seed of C and D, supply the same question and successive prefixes of that saved answer to the base and tuned checkpoints. Record how often they predict the same most likely next token and how far their full next-token distributions differ in each KL direction. This measures changes along answers the base could already produce. The benchmark scores still determine utility preservation.

Reuse SecPO's causal shift and response mask to align next-token logits with scored answer positions, including the registered end-of-response token. Use the pinned OP2SD \href{https://github.com/MBZUAI-reasoninglab/OP2SD/blob/dee6efb2169dfff9b068ea8cda6268182af22a2b/src/opsd_trainer.py#L1226-L1340}{full-vocabulary divergence function} with the base as teacher and the tuned model as student, in evaluation mode with the same tokenizer. Pass aligned logits in FP32 or higher, temperature one and no vocabulary truncation, clipping or sample weights. Its interpolation settings zero and one give base-to-tuned and tuned-to-base KL, respectively. Sum over scored positions across batches, then divide once by the total number of those positions. Use the same denominator for top-token agreement. On an unequal-length fixture, check both directions against direct calculation and verify that comparing the base with itself gives zero divergence and full agreement. The paper supplies the selection and scoring protocol. This experiment freezes its own item manifest.

\textbf{Qwen deployment comparison.} Section~4.6 and Table~8 test whether the defense survives a change in inference precision and decoding \cite{chen2026secpo}. Evaluate the Qwen base and optimized-data CRPO checkpoint in BF16, in llama.cpp Q8\_0, and in Q8\_0 with the same pretrained DFlash drafter. Pin the public \href{https://huggingface.co/z-lab/Qwen3.6-27B-DFlash/tree/0919688658996800f86b895034249700e9481106}{Qwen3.6-27B DFlash drafter} for CRPO. For llama.cpp, use the same \href{https://huggingface.co/ggml-org/Qwen3.6-27B-GGUF/blob/ce573e705fbf0d2cd70a24799c065ef72b50337c/dflash-Qwen3.6-27B-Q8_0.gguf}{ggml-org Q8\_0 drafter file} with both targets. Verify its 1,849,481,696 bytes against SHA-256 {\small\texttt{120f\allowbreak{}9f2d\allowbreak{}b399\allowbreak{}1a03\allowbreak{}6c1c\allowbreak{}18ee\allowbreak{}4441\allowbreak{}e1bc\allowbreak{}6165\allowbreak{}569f\allowbreak{}7c7f\allowbreak{}9512\allowbreak{}d8ff\allowbreak{}ed8c\allowbreak{}930d\allowbreak{}e687}} before inference. Its GGUF v3 header identifies architecture \texttt{dflash}. This conversion was published on November~8, 2026, after the paper. SecPO does not identify its historical drafter file by hash. Use llama.cpp's \href{https://github.com/ggml-org/llama.cpp/blob/master/docs/speculative.md}{\texttt{draft-dflash} mode} with target verification. Pin the draft limit and report proposed and accepted draft tokens. Report DFlash throughput only when at least one draft token was proposed. Convert each merged BF16 checkpoint to Q8\_0 and record the conversion revision, weight hash, and llama.cpp build. The authors' \href{https://huggingface.co/Sizhe-Chen/Qwen3.6-27B-Instruct-SecPO-8Bit/blob/8674c489dda144c51441b5a75cfb7cfd2f7d49f8/README.md}{Hub 8-bit release} contains bitsandbytes INT8 weights and a validation record for a separate local Q8\_0 conversion. Before benchmarking, reuse its \href{https://huggingface.co/Sizhe-Chen/Qwen3.6-27B-Instruct-SecPO-8Bit/blob/8674c489dda144c51441b5a75cfb7cfd2f7d49f8/smoke_cases.json}{arithmetic and factual-question fixtures} with their published expected strings in each serving mode. Log the rendered prompt, actual template hash, stop rules, raw answer and expected-substring decision. The archived \href{https://huggingface.co/Sizhe-Chen/Qwen3.6-27B-Instruct-SecPO-8Bit/blob/8674c489dda144c51441b5a75cfb7cfd2f7d49f8/ollama_validation_results.json}{Ollama serving record} identifies its template by SHA-256 {\small\texttt{fc7c\allowbreak{}6fc0\allowbreak{}19f4\allowbreak{}a1a7\allowbreak{}5a9d\allowbreak{}a22d\allowbreak{}95dd\allowbreak{}e705\allowbreak{}2570\allowbreak{}30f8\allowbreak{}0c56\allowbreak{}6ed4\allowbreak{}6821\allowbreak{}17ce\allowbreak{}eca0\allowbreak{}a24c}}. Fingerprint the CRPO template separately. Reserve the two fixtures with an untrusted \texttt{input} role for the native SecPO packaging control. Report these checks separately from benchmark scores. Use the same fixed framed targets, benign controls, and clean utility items across modes. For reference-based utility, use the same registered comparator and judge across modes. Hold prompt tokens, greedy generation, stop rules, and output limits fixed. Report paired changes in unsafe answers, over-refusal, and utility, alongside output tokens per second on identical hardware and workloads. Compare the two Q8\_0 runs to measure the effect of the drafter. Record any runtime change in the BF16 comparison. Keep deployment results separate from the primary training comparison.

\textbf{Joint success criterion.} Before training, lock the primary Llama model comparison, endpoint definitions, benchmark versions, margins, and checkpoint rule using development data. Use the final scheduled update for the primary offline runs. Select any alternatives on validation families and report them separately. Compare optimized-data CRPO (D) with optimized-data DPO (C) for the effect of the complete reference prompt. To attribute a robustness benefit to including the source text, also require the simultaneous one-sided upper confidence bound on ASR(D) minus ASR(empty source) to fall below zero on the prespecified TAP endpoint. Use the three matched seed differences, the same held-out goals and the same treatment of unresolved outcomes as for D minus C. Include this contrast in the registered multiplicity adjustment. This tests adding the supplied text, including its extra length. A missing matched run leaves the three-seed attribution claim unresolved. Compare D with the original Llama base for preservation. Register tolerances in percentage points for the six clean text utility scores, useful completion on matched benign targets, over-refusal on those targets, and unsafe assistance on direct prohibited prompts. Tool, vision, deployment and cross-family outcomes have separate secondary interpretations. Choose margins from the largest scientifically acceptable loss, then use development-set discordance to estimate whether the planned test can resolve them. Keep an underpowered endpoint inconclusive. Orient utility differences so positive is better. Use the paired sampling units defined below, retaining all rows for each sampled clean SEP prompt. Account for repeated seeds and preregister simultaneous bounds for the primary endpoints. Claim utility preservation only when each lower bound exceeds its negative tolerance. For over-refusal and direct unsafe assistance, the upper bound on the increase must be below its tolerance. An interval touching or crossing a boundary is inconclusive. The confirmatory effect for each primary outcome is the equally weighted mean of the three seed-specific paired differences on the same test items, conditional on those fitted checkpoints. Report each difference as well. In each resample draw the common test units once for all arms and seeds, compute the three paired differences, then average them. A reused base answer stays one shared observation. For framing tests use source families, combining all families connected by a shared scored answer. Retain the registered analysis weights. For utility, preserve each benchmark's aggregation weights and use its independent item or cluster as the resampling unit. A complete, fully judged benchmark has an observed score for the saved answers. An item interval concerns generalization to the named population of comparable items under its stated sampling assumptions, conditional on the three fitted models. Three seeds give limited precision about variation across future training runs. For binary pairs sampled independently from a common population, let $p_{10}$ and $p_{01}$ be the probabilities of improvement and regression, respectively, after orienting the outcome so higher is better. Report the observed counts of improvements and regressions alongside the estimated paired difference $p_{10}-p_{01}$. For each direction, divide its count by the total number of eligible paired units, including those on which both models agree. Use the same exact routine for the component bounds. At allocated tail error $a$, use \texttt{alternative='greater'} for a lower bound and \texttt{alternative='less'} for an upper bound, each at confidence level $1-a$. Subtract opposite endpoints to bound $p_{10}-p_{01}$. For fixed stratum counts, construct simultaneous bounds within strata where the binomial model applies and combine them with the registered nonnegative population weights summing to one. Independent observations alone can have unequal discordance probabilities, in which case a pooled one-sided binomial bound can undercover (\href{https://www.math.uni.wroc.pl/~pms/files/35.2/Article/35.2.8.pdf}{Mattner and Tasto, 2015}). Allocate coverage jointly across primary endpoints, contrasts, seeds, strata and the component bounds used. For this exact binary procedure, average the three simultaneous seed-specific lower bounds and the three upper bounds to bound the registered mean. Even when there are zero observed discordances, this procedure retains finite-sample uncertainty. For clustered or weighted outcomes, preserve the registered design and treat a degenerate resampling interval as inconclusive unless a prespecified bound valid for that design resolves the margin. Report raw results, even when a suite lacks power.

Apply the primary decision rules over all completions consistent with the recorded outcomes. For each fully executed goal and budget, a confirmed attempt fixes its success indicator at one. Its upper bound is one only if some eligible attempt is confirmed or could still reach nine YES votes. Any recorded valid NO rules out unanimity on that attempt. An unresolved validity decision can contribute to the upper bound only if a valid target and a unanimous verdict are both still possible. Keep incomplete execution separate and allow unobserved attempts only within the registered budget. Bound excess failures by assigning admissible outcomes adversely across arms, while honoring every recorded vote. Apply the same adverse-direction principle to missing scores on each preservation endpoint within its scoring range. Preserve paired units, shared observations and registered weights when forming bounds. If the confidence bound for any feasible assignment reaches or crosses the decision boundary, report the confirmed scores and an inconclusive comparison. These bounds concern the registered panel and its observed search histories. Human disagreement with a complete verdict is a separate measurement error.

A primary robustness claim requires a lower confirmed ASR for D than C on the prespecified adaptive endpoint (TAP at 800), with the upper confidence bound on ASR(D) minus ASR(C) below zero, alongside the preservation criteria. A separate claim that CRPO improves robustness over the unmodified base requires the same upper-bound test for D minus base. Estimate matched benign helpfulness over all validated benign targets, including empty answers and refusals. Before viewing answers, specify what useful progress requires for each control. If a necessary fact is missing, a reply can earn credit by asking for that fact and providing any help already possible under the item's criterion. For the import example, asking which country the user will enter can earn credit. Asking for a destination the user already supplied earns no clarification credit. Use the same criterion for every model. A successfully completed generation with no visible answer earns zero helpfulness credit. Record a serving failure as missing, retry the same checkpoint and item, and retain unresolved items in coverage and paired sensitivity bounds. Report how many completed answers were empty separately from refusals. Report Genetic, learned-attacker, partial-overlap, and transfer results separately. A lower attack-success rate alone can arise from blanket refusals on mixed requests. For a claim that CRPO improves safe and useful handling of these requests, report the paired D-minus-C change in the secondary safe-and-helpful rate on the common eligible targets from both overlap classes. Keep families linked by a shared answer together and test sensitivity to unresolved outcomes. Claim improvement on this secondary outcome only if the lower confidence bound stays above zero under assignments consistent with the unresolved labels. For descriptive comparability with SecPO's Tables~3--4, also report the unweighted mean of the six clean text utility scores and the mean of Genetic ASR@800, TAP ASR@800 and PISmith ASR@10, with all component scores adjacent on a common 0--100 scale. Label the latter as a mean of three attack-specific rates with different budgets. Keep the confidence-bound tests on the prespecified individual outcomes.

Reuse Figure~1's two-axis display and Appendix~B.6's \href{https://arxiv.org/html/2610.41998v1#A2.SS6}{per-coordinate provenance table} \cite{chen2026secpo}. For each checkpoint with both measurements, plot task-weighted clean AgentDojo completion across the 97 tasks on the horizontal axis and verified framing TAP ASR at 800 defender queries on the vertical axis. Record each endpoint's numerator, denominator, interval and run source alongside the model revision and reasoning mode. The \href{https://arxiv.org/src/2610.41998v1}{SecPO source archive} contains the original Figure~1 PDF and Table~11 LaTeX. Compute the new coordinates from the registered CRPO runs. Use separate panels by family. SecPO's original vertical coordinate measures agentic injection, so keep its published coordinates in their original comparison.

\textbf{Resource accounting and failure analysis.} TAP and Genetic together allow 1,638,400 defender queries per checkpoint before early stopping. Include validity calls, nine-vote verification, retries, preference sampling, and human audits in cost estimates. Use pilot throughput and acceptance yield to project GPU time and serving cost before reserving resources. Draw 1,024 candidate records from adaptive evaluation for manual audit, or review the full population if fewer are available. Cover confirmed deliveries, unconfirmed answers, split votes and invalid reframings where observed. Assign every record to one stratum before drawing the sample. Allocate at least one draw to each nonempty stratum in the audit population, combining strata before sampling if needed. Sample uniformly without replacement within each stratum. For a stratum of $N_h$ records with $n_h$ draws, give each draw weight $N_h/n_h$.

Following SecPO's manual check of its SEP success rule \cite[Sec.~4.1]{chen2026secpo}, review target validity first. For valid targets with a saved answer, a completed nine-vote panel verdict and a determinate blinded human judgment, use scikit-learn's \href{https://scikit-learn.org/1.6/modules/generated/sklearn.metrics.confusion_matrix.html}{weighted confusion matrix}. The blinded human reviewer receives the same criterion wording as the panel and scores whether the saved answer delivered that capability. Record other policy violations separately. Encode human nondelivery and panel-unconfirmed as 0, and human delivery and panel-confirmed as 1. Pass human codes as \texttt{y\_true}, panel codes as \texttt{y\_pred}, and set \texttt{labels=[0,1]}, \texttt{sample\_weight} to the record weights, and \texttt{normalize=None}. The four cells must sum to the weight of the included records. For the false-positive rate, divide by the total weight classified by humans as nondelivery. For the false-negative rate, divide by the total weight classified by humans as delivery. Show each weighted numerator and denominator. A zero denominator makes that rate unestimable. Report the counts and represented weights of excluded cases by reason, including invalid targets, absent answers and unresolved judgments. Compare the panel with the prespecified first vote on the same cases with both readings available and report that coverage. Extra cases selected for illustration do not enter the estimates. Following Appendix~D, illustrate audited gains and regressions where observed by pairing DPO and CRPO answers to the same fixed target ID. Include its benign counterpart when examining over-refusal. Show the target wording that supports each judgment and the frequency of the outcome in the full test.

\section{Methodological differences from SecPO}
SecPO's training recipe and released offline data provide the comparison below \cite[Secs.~3--4 and App.~B]{chen2026secpo}. The primary changes arise because reframing can change which assistance is appropriate.
\begin{center}
\small
\begin{tabular}{@{}p{0.16\linewidth}p{0.36\linewidth}p{0.39\linewidth}@{}}
\hline
 & \textbf{SecPO} & \textbf{Benign framing experiment} \\
\hline
Task and reference & A trusted task remains after removing an injection from external data. The frozen model sees the clean input. & The framed request is the user's task. The reference sees both requests and answers the target under its applicable policy. \\[3pt]
Search input & A candidate instruction is inserted into the data seen by the defender. PAIR refines it from the answer and critic feedback. & The same search roles exchange a whole target request, its ordinary answer, a validity decision, a score, and a critique. The source is visible to the proposer. \\[3pt]
Training record & \texttt{prompt} holds the attacked input and \texttt{clean\_prompt} the clean task. Answers come from the trusted and injected tasks. The archived Qwen3.6 data keep the original labels. Qwen3.8 regenerates labels from earlier task inputs. & Policy input $x_i^F$, reference input $c_i$, chosen $y_i^C$, rejected $y_i^0$. Keep a row after the blinded preference check. \\[3pt]
Judgment & The search critic sees an injection goal and a defender answer. SEP uses a witness check and semantic votes on instruction following. & The policy judge reads the target alone, the relation judge reads source and target, and the answer judge reads target and answer(s). Each sees the relevant rubric. \\[3pt]
Training schedule & Online Llama search and offline Qwen transfer. & Offline factorial comparisons and matched online DPO/CRPO on Llama. Qwen3.6 crosses both objectives with both datasets. Qwen3.8 compares objectives on transferred selected targets. \\[3pt]
Evaluation & Static and adaptive attacks, agentic and visual security, and utility. & Six static frame styles, held-out adaptive text attacks, six clean text utility suites, clean AgentDojo, and clean Qwen AgentDyn and DocVQA. \\
\hline
\end{tabular}
\end{center}

\subsection{Implementation boundaries and deliverables}
\textbf{Direct reuse.} Pin the SecPO repository revision in the experiment manifest \cite{secpocode}. Use the training-time PAIR worker for full construction, the portable search fixtures for replay, and the TAP and Genetic controllers for evaluation. Reuse \href{https://github.com/Sizhe-Chen/SecPO/blob/da0372a0e598badf4e51c98c175a962147cac2ac/src/secpo/loss.py}{masked sequence scoring}, adapter-disabled reference passes, and the sigmoid loss. Record changes to callbacks, data validation, and budget settings against that revision. Retain the original injection fixtures as regression checks.

\textbf{Independent ordinary-DPO check.} Use two development preference pairs with unequal answer lengths and a small model whose policy and frozen reference give different scores. The local TRL launcher pools validation and test rows for evaluation during training. Use a scoring-only call on the fixture. CRPO's exported answers already contain a terminal token. TRL~0.26.0's \href{https://github.com/huggingface/trl/blob/84794a7d03c748f3e1518cd3efc9a944b8cdd2bd/trl/trainer/dpo_trainer.py#L670-L735}{text tokenizer} appends another to each answer, which can change the loss. Pass the ordinary prompt and each unpadded response as token IDs to TRL's preference collator, retaining the single scored terminal token. With dropout off, match the four sequence scores to SecPO on the same weights. Use TRL's \href{https://github.com/huggingface/trl/blob/84794a7d03c748f3e1518cd3efc9a944b8cdd2bd/trl/trainer/dpo_trainer.py#L1008-L1102}{sigmoid loss} with the same $\beta$, reverse-KL divergence, reference scoring enabled and zero label smoothing. Check each loss and its derivatives with respect to the two student scores, then exchange the answer labels and repeat.

\textbf{Online exchange and snapshots.} Adapt the typed request and result schema in \href{https://github.com/Sizhe-Chen/SecPO/blob/da0372a0e598badf4e51c98c175a962147cac2ac/src/secpo/online_training.py}{\texttt{online\_training.py}} as well as the dataset loader and collator. The exchange also checks for deletion of a contiguous injection. The upstream worker returns only an attacked-prompt string per ID, and its collator replaces only the policy prompt. Version the exchange to return a complete accepted record or a skipped outcome. Carry source and target IDs, both rendered contexts, both fresh answers, validity, preference decision, and serving snapshot. Collate these fields atomically. Persist accepted record IDs, reserve cursors and sampler state with the update so a retry replays the same batch. Use the worker's per-round and final snapshot lookup, readiness checks, ordered batching, and per-update adapter publication. Persist \texttt{last\_results} after each request and resume from a consistent checkpoint and dataset cursor. Apply the same freshness policy to online DPO and CRPO. The public synchronous microbatch worker can see the same snapshot throughout a search. For the full online runs, use the external producer required by the strict worker and pipeline the next group while the trainer consumes the previous one. Each worker request contains 32 distinct, eligible sources for harmful-framing search per rank. After search and blinded preference judging, the producer assembles the optimizer batch. At the planned equal mixture, each rank takes the first 16 accepted harmful pairs in prespecified source order and 16 accepted benign controls from a shared, independently constructed training queue. Across eight ranks this gives 128 pairs of each class per update. Controls bypass PAIR. For a shortfall, use full reserve search blocks under a common prespecified per-update budget and source schedule for both online arms. Log all attempted, rejected, surplus, and consumed rows. If the reserves cannot fill the quotas, flag the update as incomplete before either arm advances. This preserves the native worker width while accounting for CRPO's preference filtering. If the synchronous path is used, record it as a different schedule. Record which schedule ran and every observed snapshot change \cite{secpocode}.

\textbf{Reproducibility checks.} Replay a deterministic batch against the pinned training backend and its adapted version, comparing stream histories, ties, empty proposals, final-generation calls, vote order, and snapshot IDs. Use the data-audit and merge tools to record input IDs, row counts, token-length exclusions, checksums, and LoRA merge provenance. For the published SecPO controls, replay the nested-delimiter examples in the \href{https://huggingface.co/Sizhe-Chen/Llama-3.1-8B-Instruct-SecPO/blob/8b019ad9b49d48911d6dcff2fc8861be0cce706d/README.md}{Llama} and \href{https://huggingface.co/Sizhe-Chen/Qwen3.6-27B-Instruct-SecPO/blob/726da6ba0f153d478c207dd5edc593e18e7fd513/README.md}{Qwen model cards} through the tokenizer used for evaluation. Use Llama's supplied token-ID equality assertion. For Qwen, check that the rendered attack text remains inside its original input message after repeated delimiter removal. Pin the distributed template and record which renderer produced the tokens. The published template hashes changed after the initial merges. The Llama merge record identifies the earlier template (\texttt{04c6256d}), and the fixed-point filtering revision has prefix \texttt{307ee507}. The Qwen record identifies the base template (\texttt{e84f32a2}), its first SecPO upload has prefix \texttt{e30baf3a}, and its fixed-point revision has prefix \texttt{fc7c6fc0}. The standalone and tokenizer-embedded copies agree at these historical revisions. In the later \href{https://huggingface.co/Sizhe-Chen/Llama-3.1-8B-Instruct-SecPO/tree/d244fe7440209ff089418ca8d0717ef7bda0da4c}{Llama SecPO release} and the matching \href{https://huggingface.co/Sizhe-Chen/Llama-3.1-8B-Instruct-MetaSecAlign/tree/5ae69e7ba7fe1ff8b5cdd4f24eb6aad1c54c6794}{Meta-SecAlign release}, each tokenizer ships a standalone template and omits the embedded copy. Record which file or benchmark wrapper actually rendered the clean tool transcript for each serving revision. Hold ordinary user-message rendering fixed across the matched DPO and CRPO arms. Before relying on the adapted evaluation path, run the \href{https://github.com/Sizhe-Chen/SecPO/blob/da0372a0e598badf4e51c98c175a962147cac2ac/src/secpo/evaluation/sep_static.py}{unchanged static SEP runner} on the canonical 9,160-item injection test and the released same-family checkpoints. Select \texttt{PAPER\_STATIC\_ATTACKS} explicitly. The CLI defaults also produce two model-specific Basic Adaptive attacks in a separate category. Apply the native witness check and the \href{https://github.com/Sizhe-Chen/SecPO/blob/da0372a0e598badf4e51c98c175a962147cac2ac/scripts/judge_sep_outputs.py}{released nine-vote judge} to the outputs. Rejoin verdicts by source and attack, check for six distinct results per source, and report the static category's any-of-six ASR. Its native denominator excludes pending sources. Record pending counts and all-source lower and upper bounds before comparing with the paper. Use the authors' \href{https://huggingface.co/Sizhe-Chen/Llama-3.1-8B-Instruct-MetaSecAlign/blob/dbc4ed4aa435b651c85d47c154a8908fe0f7129f/evaluation_results.json}{machine-readable Table~1} as the historical Llama reference. Its 30 entries across ten metrics and three checkpoints match the released \href{https://github.com/Sizhe-Chen/SecPO/blob/da0372a0e598badf4e51c98c175a962147cac2ac/docs/PERFORMANCE.md}{performance table}. The authors identify these as manuscript values transcribed at export. For the six static attacks it reports 98.1\% for the base, 1.8\% for Meta-SecAlign and 1.4\% for SecPO. Record fresh measurements in separate columns. Record dataset, template, checkpoint, and judge revisions. Investigate discrepancies before attributing a change in ASR to CRPO. Keep this original-task control in a separate results table. Run the native missing-vote, dense-score, reasoning-removal and resume fixtures, then exercise the corresponding semantic-validity cases. In the framing tests, a harmless answer that quotes the prohibited request must remain eligible for judging and receive a negative violation verdict. For utility, verify the identity and hash of each external reference file and the exact order of the scored prompts. First run the unchanged upstream loss, collator, and judge fixtures. On identical injection records, check all four response log probabilities and the per-example loss in the adapted trainer against upstream. Exercise invalid rewrites, split votes, an exhausted vote, and resume on framing fixtures. Archive evaluator settings and raw outcomes with checkpoint IDs. Report paper settings, released defaults, and overrides distinctly whenever they differ. 

\textbf{Native preference-data controls.} Reconstruct the five parts of the archive using its \href{https://github.com/Sizhe-Chen/SecPO/blob/da0372a0e598badf4e51c98c175a962147cac2ac/data/qwen36_offline_dataset/manifest.json}{manifest}. Audit the raw rows, run the original preference preparer, then audit the four-field training rows. At the pinned revision, both local audits counted 19,157 rows and zero duplicate complete rows. The \href{https://github.com/Sizhe-Chen/SecPO/blob/da0372a0e598badf4e51c98c175a962147cac2ac/data/DATASET_CARD.md}{released dataset card} documents 19,145 rows after the separate Qwen prompt-length filter \cite{secpocode}. Preserve the original bytes and use separate manifests for native SecPO records and constructed CRPO pairs. The \href{https://huggingface.co/datasets/Sizhe-Chen/Qwen3.6-27B-Instruct-SecPO-trainset/blob/0cbfd6e6a2c2966dbccf4fa89b6ea7fe20c8703e/README.md}{public Qwen dataset} and the reconstructed archive both contain 99,314,808 raw bytes with SHA-256 \nolinkurl{cd349baa406eae0bf7e4ee2d938d087d6656469d88076d68ba76baeb7e5ec669}. For the Llama input-control fixture, use the \href{https://huggingface.co/datasets/Sizhe-Chen/Llama-3.1-8B-Instruct-SecPO-trainset/blob/f347c3abe1f451e88b02918c6bc3921e6816d8c6/README.md}{released optimized Llama snapshot}. Its 81,546,526-byte JSON file matched SHA-256 \nolinkurl{899f073891a6e1579c723a93ae0e838ea2b24eba634baa8332504e2c7d007f73}. In \texttt{--online-pair} mode, the native preparer produced 19,157 unique IDs and 19,157 nonempty goals. It supplies final optimized inputs for offline reuse; recovering the historical online trajectory would additionally require the ordered worker results and checkpoint history. Check the normalized attack goal, stable IDs on replay, and rejection of duplicate IDs \cite{secpocode}. For CRPO, assign the immutable source-family ID at the split and a separate search-example ID before sampling answers. Key the starting styles and replay to these IDs. Fingerprint the completed preference pair separately. The native fallback ID includes both answers, so regenerating an answer would otherwise change which search is replayed.

\textbf{Scope of parity.} Preserve the upstream checkpoint identity, native chat template and distinct search, verifier, and trainer renderings in the injection reproduction control. In native PAIR, search and final verification append the candidate to the data with different separator handling, and training restores the original insertion slot. In CRPO, all three stages refer to a complete standalone target. Require the target text and ordinary deployment rendering to agree across those stages and record their hashes separately from the combined reference. The public trainer refactors the original Qwen3.6 four-pass implementation. For Llama it replaces the historical Torchtune training and exchange path with a Hugging Face implementation. Unretained training randomness and historical hosted-model versions limit comparisons to protocol and measured outcomes. Record actual model routes, seeds, scheduling and data-selection differences with each result \cite{secpocode}.

\textbf{What the project must supply.} Across the releases considered here, the remaining work is the policy-specific content and its validation. This includes the combined reference instruction, the criteria for capability overlap and benign controls, calibrated judgments and adjudicated examples, and the links between each source, rewrite, control, and fresh pair of answers. Existing model-serving, red-teaming, and preference-training tools can execute these operations once their inputs and decisions are specified. WildJailbreak supplies candidate seeds and styles. XSTest supplies the external over-refusal test, and the published WildJailbreak evaluation provides a secondary fixed framing test. OR-Bench is a related over-refusal benchmark \cite{jiang2024wildteaming,rottger2023xstest,cui2024orbench}. On domains with an exact behavior match, HarmBench's existing completion classifier can provide an additional check on whether an answer delivers that behavior \cite{harmbench}. On saved human-audit answers, the released \href{https://github.com/allenai/wildguard/blob/495b9dc837be594c6c5665b0f7f56169baf0d5be/wildguard/wildguard.py#L123-L167}{WildGuard evaluator and parser} also provide separate refusal and response-harm judgments. Compare successfully parsed binary fields with their corresponding audited labels. Record the row-level \texttt{is\_parsing\_error} flag and field-level missingness separately. Exclude flagged rows from agreement estimates and report their count even when a field contains a value. On full-overlap targets with a clearly forbidden standalone request, the \href{https://github.com/dsbowen/strong_reject/blob/7a551d5b440ec7b75d4f6f5bb7c1719965b76b47/strong_reject/evaluate.py#L378-L427}{StrongREJECT rubric} can independently score how useful the saved response is for carrying out that request. The score is gated by refusal. Give the \href{https://github.com/dsbowen/strong_reject/blob/7a551d5b440ec7b75d4f6f5bb7c1719965b76b47/strong_reject/evaluate.py#L181-L201}{judge helper} a single registered model route for all arms. Its model field records the requested route. Its \href{https://github.com/dsbowen/strong_reject/blob/7a551d5b440ec7b75d4f6f5bb7c1719965b76b47/strong_reject/generate.py#L69-L86}{generation helper} returns only message text. Save the raw reply and any provider-reported model identity at the serving boundary. Record the observed identity as unverified when the provider omits it. Require one binary refusal rating and one integer from 1 to 5 each for convincingness and specificity. Report missing, ambiguous, out-of-range or nonfinite ratings and confirmed route mismatches as unscored, with counts and coverage. Use these as secondary checks of saved outcomes within their respective scopes. The policy-aware assessment determines partial refusal and useful permitted assistance. A source behavior label alone cannot decide partial overlap or permitted help under the framed request. Those decisions require project-specific calibration.

\textbf{Integration milestones.} Document model revisions, instrument search and verification, validate four-way scoring, then run the dedicated feasibility pilot. Freeze the full protocol before training. Link the utility harnesses to the same scored checkpoint manifest. The local prototype supplies model-call and replay backends, JSONL audit records, and a standard-DPO script.

\section{Related work and proposed contribution}
SecAlign applies preference optimization to prompt-injection defense; SecPO changes the reference context to the corresponding clean prompt \cite{chen2024secalign,chen2026secpo}. This proposal retains SecPO's loss structure while constructing a reference for benign framing, where an unchanged source answer may be inappropriate.

Synthetic alignment data and prompted self-improvement are established. Constitutional AI uses policy-guided model feedback; Counterfactual DPO constructs alignment preferences through prompting; RLCD generates contrasting responses with different prompts \cite{bai2022constitutional,butcher2024counterfactual,yang2023rlcd}. WildJailbreak includes harmful and benign adversarial examples \cite{jiang2024wildteaming}. TeachBack also samples answers from one model with and without privileged context. Its reasoning experiment trains in both directions with reversed pair preferences. CRPO uses the target-only safety judgment to select one direction and places the combined context in the frozen reference term \cite{ji2026teachback}. Offset DPO already adjusts the required preference margin per example \cite{amini2024odpo}. Intent-Privilege OPSD uses training-only evidence about intent to distill safe, useful responses into a vision-language model \cite{deng2026intent}. SecPO itself describes suitable reference distributions beyond injection as a research direction \cite[Secs.~6--7]{chen2026secpo}. In text safety, OPSA measures how often a privileged context changes an unsafe answer into a safe one before applying on-policy distillation \cite{fu2026opsa}. Safe-Completions uses policy-conditioned generation in its supervised stage and assesses the usefulness of assistance within the safety boundary \cite{yuan2025safecompletions}. OR-Bench already rewrites harmful seeds into benign requests and validates the results \cite[Sec.~3.2]{cui2024orbench}. OpenSafeIntent evaluates matched benign, dual-use and malicious versions of a task and measures safety changes across paraphrases \cite{uppaal2026opensafeintent}. RiskCueBench directly elicits the evaluated model's interpretation and pairs it with that model's response \cite{wang2026riskcue}. CRPO's target-only assessor establishes which reading the visible text supports; the experiment measures the trainee's answer and leaves its recognition of intent unmeasured. GoalFrameBench independently varies goals and framings to study a detector on a frozen model \cite{farzam2026goalframe}. Its preprint describes the design and examples and states that it omits the full corpus.

Irpan et al.'s DPO baseline already generates the preferred answer to a direct request and the rejected answer to its wrapped version, then trains on the wrapped request \cite[Sec.~3.2]{irpan2025consistency}. Han et al.'s on-policy consistency training scores the student's response to a rewritten request token by token under the direct request \cite[Sec.~4.2]{han2026opct}. Both designs address transformations that preserve the requested behavior. CRPO also studies transformations that change which help is appropriate. Its combined reference has to answer the visible target correctly even when the source alone warrants a different response.

The rDPO baseline in daDPO uses teacher-generated preferred answers, student-generated rejected answers and the teacher distribution as the reference for both answers on the same input \cite[Sec.~4.1]{zhang2025dadpo}. CRPO reuses the SecPO loss and one checkpoint under two contexts. A target-only judgment determines whether the candidate pair enters training. The empty-source training control estimates what adding the source text contributes beyond the rubric. The development scoring check asks whether the identity of the source matters for the reference margin.

E-DPO established a safety application of cross-context DPO in 2024. Its student scores responses on the original prompt while the frozen reference can score them on a safety-transformed prompt \cite[App.~C]{su2024edpo}. CRPO's combined reference includes an earlier request that can share all, part, or none of the target's prohibited capability. Its matched benign examples test what happens when that earlier request would justify a restriction the visible target has removed.

SHARD also constructs safety-training data by comparing a direct answer with a rewritten one. The model receives the original request, a safer reframing and its direct answer. A judge selects rewrites that improve helpfulness, and supervised fine-tuning pairs those answers with the original request \cite[Secs.~3.2--3.4]{manoranjan2026shard}. The authors' release also \href{https://github.com/Viswonathan06/shard-self-reframing/blob/dcf5eba32138951732e5d5d0590c44b1e18e9e40/scripts/prepare_sft_dpo_data.py#L73-L85}{exports the saved winner and loser under the original request} for an \href{https://github.com/Viswonathan06/shard-self-reframing/blob/dcf5eba32138951732e5d5d0590c44b1e18e9e40/scripts/train_dpo_linguasafe.py#L142-L150}{ordinary DPO trainer}. The paper's reported HH-Tune baseline uses Anthropic HH preference data \cite[Sec.~4.3]{manoranjan2026shard}. CRPO makes the framed request the deployment input. The policy judgment therefore turns on what that request still asks for, including cases where only part of the original restriction survives.

PBSD also pairs a teacher-context answer with an ordinary student answer. It samples the latter from the current student, and its released trainer can filter pairs using known task answers \cite{yu2026pbsd}. In CRPO's primary comparison, the frozen checkpoint generates both answers before training. A judge seeing only the target, rubric and candidate answers must prefer the combined-context answer in both presentation orders. This tests whether the trained model withholds prohibited assistance on one target and gives useful help on its matched benign counterpart.

The proposed contribution is the construction and evaluation of a combined source-and-target reference that answers the target under its applicable policy, with counterfactuals testing inappropriate transfer of the source's harmful label. Novelty remains a qualified application-level claim, subject to broader literature comparison and experimental results. The project does not claim that DPO has previously been used only for prompt injection or that this loss is mathematically new.

\section{Remaining uncertainty and practical scope}
Full-study deliverables include the offline datasets, both online search histories, Qwen transfer and ablation records, primary and diagnostic training comparisons, adaptive robustness curves, over-refusal metrics, and a per-checkpoint utility scorecard. Release configuration, dataset attributions and aggregate audit results. For source text carried into combined prompts, document the applicable access and redistribution terms. Where raw text is distributed through its upstream source, provide the revision, row locator, digest and rendering instructions so authorized readers can reconstruct the same record. Run SecPO's \href{https://github.com/Sizhe-Chen/SecPO/blob/da0372a0e598badf4e51c98c175a962147cac2ac/scripts/validate_release_manifest.py}{source-manifest validator} on a clean extraction of the CRPO source archive using its own exact file inventory. On a clean archive of upstream \texttt{da0372a}, the original validator flags the README email rule, the bundled paper PDF and the unlisted AgentDojo example JSON. Record that historical result separately. Report packaging results separately from the dataset checksums and benchmark measurements. The scope is semantic framing of text requests. Image injection, white-box token attacks, and attacks delivered through untrusted tool output address different threat models. Appendix C.3 also optimizes continuous attack embeddings. A user submitting a natural-language framing controls the text, so this study adopts the paper's discrete text search. Clean tool use and Qwen document-image understanding measure utility transfer.

The central uncertainty is empirical: the combined context must supply useful training signals without importing unsupported assumptions or excessive refusal. The pilot tests this premise; a well-defined objective cannot establish it by itself.

Development will use approved existing safety data or harmless fictional policy tasks, following course requirements. Sampling on safety prompts may produce unsafe outputs; retained inferior outputs are training negatives, not desired behavior. The goal is to improve safety while preserving benign helpfulness.

\begin{thebibliography}{99}

\bibitem{irpan2025consistency}
Alex Irpan, Alexander Matt Turner, Mark Kurzeja, David K. Elson, and Rohin Shah.
\newblock \href{https://arxiv.org/abs/2510.27062}{Consistency Training Helps Stop Sycophancy and Jailbreaks}. arXiv:2510.27062, 2025.

\bibitem{han2026opct}
Andy Han, Kristina Fujimoto, Avidan Shah, et al.
\newblock \href{https://arxiv.org/html/2605.21834v1}{On-Policy Consistency Training Improves LLM Safety with Minimal Capability Degradation}. arXiv:2605.21834, 2026.

\bibitem{dhouib2026raised}
Mohamed Dhouib, Clement Elliker, Alexi Canesse, et al.
\newblock \href{https://arxiv.org/html/2610.06401v2#A2.SS1}{RAISED: Self-Distillation for Robustness to Prompt Injection in LLM Agents}. arXiv:2610.06401v2, 2026.

\bibitem{davion2026drift}
Ravenor Davion and Nick Rui.
\newblock \href{https://arxiv.org/html/2610.07842v1#S2.SS4}{Privileged Context as Drift in On-Policy Self-Distillation}. arXiv:2610.07842, 2026.

\bibitem{rafailov2023dpo}
Rafael Rafailov, Archit Sharma, Eric Mitchell, Christopher D. Manning, Stefano Ermon, and Chelsea Finn.
\newblock Direct Preference Optimization: Your Language Model Is Secretly a Reward Model.
\newblock Advances in Neural Information Processing Systems, 2023.

\bibitem{zheng2023mtbench}
Lianmin Zheng, Wei-Lin Chiang, Ying Sheng, et al.
\newblock \href{https://arxiv.org/abs/2306.05685}{Judging LLM-as-a-Judge with MT-Bench and Chatbot Arena}.
\newblock Advances in Neural Information Processing Systems, 2023.

\bibitem{uppaal2026opensafeintent}
Rheeya Uppaal, Seungwoo Lyu, Selina Sung, and Junjie Hu.
\newblock \href{https://arxiv.org/abs/2607.02047}{OpenSafeIntent: Evaluating Intent-Calibrated Safe Completion Across Dual-Use Prompt Sets}.
\newblock arXiv:2607.02047, 2026.

\bibitem{wang2026riskcue}
Yuwei Wang, Xingwang Liang, Yu Xie, Haotian Chi, and Haijun Geng.
\newblock \href{https://arxiv.org/abs/2610.22425}{Recoverability-Controlled Safety Miscalibration in Large Language Models}.
\newblock arXiv:2610.22425, 2026.

\bibitem{chen2024secalign}
Sizhe Chen, Arman Zharmagambetov, Saeed Mahloujifar, Kamalika Chaudhuri, David Wagner, and Chuan Guo.
\newblock SecAlign: Defending Against Prompt Injection with Preference Optimization.
\newblock arXiv:2410.05451, 2024.

\bibitem{chen2026secpo}
Sizhe Chen, Yibo Peng, Jaewon Chang, Chawin Sitawarin, and David Wagner.
\newblock SecPO: Principled Adversarial Training for Prompt Injection Security.
\newblock \href{https://arxiv.org/html/2610.41998v1}{arXiv:2610.41998}, 2026.

\bibitem{secpocode}
Sizhe Chen et al. SecPO public implementation and reproducibility notes.
\newblock Inspected revision \href{https://github.com/Sizhe-Chen/SecPO/tree/da0372a0e598badf4e51c98c175a962147cac2ac}{\texttt{da0372a}}.
\newblock \href{https://github.com/Sizhe-Chen/SecPO/tree/da0372a0e598badf4e51c98c175a962147cac2ac/src/secpo/attacks}{Search drivers},
\href{https://github.com/Sizhe-Chen/SecPO/blob/da0372a0e598badf4e51c98c175a962147cac2ac/src/secpo/data.py}{data validation and collation},
\href{https://github.com/Sizhe-Chen/SecPO/blob/da0372a0e598badf4e51c98c175a962147cac2ac/src/secpo/trainer.py}{trainer}, and
\href{https://github.com/Sizhe-Chen/SecPO/blob/da0372a0e598badf4e51c98c175a962147cac2ac/REPRODUCIBILITY.md}{reproducibility notes}.

\bibitem{chen2026repeat}
Sizhe Chen, Yu-Lin Tsai, Ivan Evtimov, et al.
\newblock \href{https://arxiv.org/abs/2609.04533}{Repeat-After-Me: Black-Box Adaptive Visual Prompt Injection}. 2026.

\bibitem{mathew2021docvqa}
Minesh Mathew, Dimosthenis Karatzas, and C.~V. Jawahar.
\newblock \href{https://openaccess.thecvf.com/content/WACV2021/html/Mathew_DocVQA_A_Dataset_for_VQA_on_Document_Images_WACV_2021_paper.html}{DocVQA: A Dataset for VQA on Document Images}. WACV, 2021.

\bibitem{peng2026secopd}
Yibo Peng, Long Lian, David Wagner, and Sizhe Chen.
\newblock SecOPD: Mitigating Adaptive Prompt Injections by On-Policy Distillation.
\newblock \href{https://aclanthology.org/2026.emnlp-main.1102/}{EMNLP, 2026}. \href{https://github.com/pppyb/SecOPD}{Released implementation}.

\bibitem{andriushchenko2025tense}
Maksym Andriushchenko and Nicolas Flammarion.
\newblock \href{https://arxiv.org/abs/2407.11969}{Does Refusal Training in LLMs Generalize to the Past Tense?} ICLR, 2025.

\bibitem{amini2024odpo}
Afra Amini, Tim Vieira, and Ryan Cotterell.
\newblock \href{https://aclanthology.org/2024.findings-acl.592/}{Direct Preference Optimization with an Offset}. Findings of ACL, 2024.

\bibitem{ji2026teachback}
Linfei Ji, Jiayun Li, Haozhao Wang, et al.
\newblock \href{https://arxiv.org/abs/2610.29167}{TeachBack: Learning to Teach via Bidirectional Preference Optimization for On-Policy Self-Distillation}. arXiv:2610.29167, 2026.

\bibitem{deng2026intent}
Haotian Deng, Wenbin Xing, Gang Xu, et al.
\newblock \href{https://arxiv.org/abs/2609.37837}{Can Vision-Language Models Stay Helpful When Facing Implicit Risks? Intent-Privilege OPSD for Efficient Safety-Helpfulness Alignment}. arXiv:2609.37837, 2026.

\bibitem{su2024edpo}
Jingtong Su, Julia Kempe, and Karen Ullrich.
\newblock \href{https://proceedings.neurips.cc/paper_files/paper/2024/hash/439bf902de1807088d8b731ca20b0777-Abstract-Conference.html}{Mission Impossible: A Statistical Perspective on Jailbreaking LLMs}. Advances in Neural Information Processing Systems, 2024.

\bibitem{manoranjan2026shard}
Viswonathan Manoranjan, Amogh Gupta, Anvesh Rao Vijjini, Thomas Hofweber, and Snigdha Chaturvedi.
\newblock \href{https://aclanthology.org/2026.findings-emnlp.1460/}{SHARD: Safe and Helpful Alignment via Self-Reframing Distillation}. Findings of the Association for Computational Linguistics: EMNLP 2026.

\bibitem{yu2026pbsd}
Xin Yu, Liuchen Liao, Yiwen Zhang, Yingchen Yu, Lingzhou Xue, and Qinzhen Guo.
\newblock \href{https://arxiv.org/html/2605.05040v2#S3.SS2}{Preference-Based Self-Distillation: Beyond KL Matching via Reward Regularization}. arXiv:2605.05040v2, 2026.

\bibitem{harmbench}
Mantas Mazeika et al. HarmBench: A Standardized Evaluation Framework for Automated Red Teaming and Robust Refusal.
\newblock \href{https://github.com/centerforaisafety/HarmBench}{Official implementation}, 2024.

\bibitem{fu2026opsa}
Yu Fu et al.
\newblock \href{https://arxiv.org/abs/2605.15239}{Reducing the Safety Tax in LLM Safety Alignment with On-Policy Self-Distillation}.
\newblock arXiv:2605.15239, 2026.

\bibitem{shrestha2026privileged}
Samyak Shrestha and Alexander Tessier.
\newblock \href{https://arxiv.org/abs/2608.18271}{Rethinking Privileged Information in On-Policy Self-Distillation}.
\newblock arXiv:2608.18271, 2026.

\bibitem{ichihara2026privileged}
Yuki Ichihara, Naoto Iwase, Mohammad Atif Quamar, and Junpei Komiyama.
\newblock \href{https://arxiv.org/abs/2608.09228}{Privileged Solutions or Context-Induced Teacher Behavior? Dissecting On-Policy Self-Distillation}.
\newblock arXiv:2608.09228, 2026.

\bibitem{zhang2025dadpo}
Zhengze Zhang et al.
\newblock \href{https://aclanthology.org/2025.findings-acl.796/}{daDPO: Distribution-Aware DPO for Distilling Conversational Abilities}.
\newblock Findings of ACL, 2025.

\bibitem{yuan2025safecompletions}
Yuan Yuan et al.
\newblock \href{https://arxiv.org/abs/2508.09224}{From Hard Refusals to Safe-Completions: Toward Output-Centric Safety Training}.
\newblock arXiv:2508.09224, 2025.

\bibitem{farzam2026goalframe}
Amirhossein Farzam, Majid Behabahani, Mani Malek, Yuriy Nevmyvaka, and Guillermo Sapiro.
\newblock \href{https://arxiv.org/abs/2602.19396}{Hiding in Plain Text: Detecting Concealed Jailbreaks via Activation Disentanglement}.
\newblock arXiv:2602.19396, 2026.

\bibitem{bai2022constitutional}
Yuntao Bai, Saurav Kadavath, Sandipan Kundu, et al.
\newblock Constitutional AI: Harmlessness from AI Feedback.
\newblock arXiv:2212.08073, 2022.

\bibitem{lopez2026boundary}
Alejo López-Ávila, Iker García-Ferrero, Jezabel Garcia, Antonio Tiene, and Román Orús.
\newblock \href{https://arxiv.org/abs/2609.04482}{Safety for Whom? Boundary-Aware Self-Distillation for Controlled LLM Safety Refusal}.
\newblock arXiv:2609.04482, 2026.

\bibitem{dubois2024length}
Yann Dubois, Balázs Galambosi, Percy Liang, and Tatsunori B. Hashimoto.
\newblock \href{https://arxiv.org/abs/2404.04475}{Length-Controlled AlpacaEval: A Simple Way to Debias Automatic Evaluators}.
\newblock arXiv:2404.04475, 2024.

\bibitem{wu2026policyalign}
Chang Wu, Junfeng Fang, Houcheng Jiang, et al.
\newblock \href{https://arxiv.org/abs/2606.25442}{PolicyAlign: Direct Policy-Based Safety Alignment for Large Language Models}.
\newblock arXiv:2606.25442, 2026.

\bibitem{butcher2024counterfactual}
Bradley Butcher.
\newblock Aligning Large Language Models with Counterfactual DPO.
\newblock arXiv:2401.09566, 2024.

\bibitem{yang2023rlcd}
Kevin Yang, Dan Klein, Asli Celikyilmaz, Nanyun Peng, and Yuandong Tian.
\newblock \href{https://proceedings.iclr.cc/paper_files/paper/2024/hash/5bd09a559a8c8e230697107b0f353d39-Abstract-Conference.html}{RLCD: Reinforcement Learning from Contrastive Distillation for LM Alignment}.
\newblock International Conference on Learning Representations, 2024.

\bibitem{jiang2024wildteaming}
Liwei Jiang, Kavel Rao, Seungju Han, et al.
\newblock WildTeaming at Scale: From In-the-Wild Jailbreaks to (Adversarially) Safer Language Models.
\newblock Advances in Neural Information Processing Systems, 2024.

\bibitem{gupta2026systemprompts}
Mohan W. Gupta and Sandy Tanwisuth.
\newblock \href{https://arxiv.org/abs/2611.11074}{System Prompts Affect Response Bias More Than Discriminability in LLM Safety}.
\newblock arXiv:2611.11074, 2026.

\bibitem{rottger2023xstest}
Paul Röttger, Hannah Rose Kirk, Bertie Vidgen, Giuseppe Attanasio, Federico Bianchi, and Dirk Hovy.
\newblock XSTest: A Test Suite for Identifying Exaggerated Safety Behaviours in Large Language Models.
\newblock Proceedings of NAACL, 2024.

\bibitem{cui2024orbench}
Justin Cui, Wei-Lin Chiang, Ion Stoica, and Cho-Jui Hsieh.
\newblock OR-Bench: An Over-Refusal Benchmark for Large Language Models.
\newblock arXiv:2405.20947, 2024.

\end{thebibliography}

\end{document}
